\documentclass[11pt]{article}

\usepackage{fontspec}
\defaultfontfeatures{Renderer=Harfbuzz,Ligatures=TeX}
\directlua{luaotfload.add_fallback("clclatin", {"name:Noto Sans:mode=harf"})}
\directlua{luaotfload.add_fallback("clclatinmono", {"name:Noto Sans Mono:mode=harf"})}
\setmainfont[Script=Arabic,RawFeature={fallback=clclatin}]{Noto Sans Arabic}
\setsansfont[Script=Arabic,RawFeature={fallback=clclatin}]{Noto Sans Arabic}
\setmonofont[Script=Arabic,RawFeature={fallback=clclatinmono}]{Noto Sans Arabic}
\usepackage[margin=1in]{geometry}
\usepackage{amsmath,amssymb}
\usepackage{booktabs}
\usepackage{array}
\usepackage{longtable}
\usepackage{tabularx}
\usepackage{graphicx}
\usepackage{caption}
\usepackage{enumitem}
\usepackage{ragged2e}
\usepackage{hyperref}
\usepackage{url}
\usepackage[bidi=basic]{babel}
\babelprovide[main,import]{arabic}
\babelprovide[import]{english}
\babelfont[english]{rm}{Noto Sans}
\newfontfamily\arabicfont[Script=Arabic,RawFeature={fallback=clclatin}]{Noto Sans Arabic}
\newfontfamily\arabicfontsf[Script=Arabic,RawFeature={fallback=clclatin}]{Noto Sans Arabic}
\newfontfamily\arabicfonttt[Script=Arabic,RawFeature={fallback=clclatinmono}]{Noto Sans Arabic}
\newfontfamily\latinfont{Noto Sans}
\newcolumntype{P}[1]{>{\raggedleft\arraybackslash}p{#1}}
\captionsetup{font=small,labelfont=bf}
\setlength{\emergencystretch}{6em}
\hypersetup{
  unicode=true,
  colorlinks=true,
  linkcolor=blue,
  citecolor=blue,
  urlcolor=blue,
  pdflang={ar},
  pdftitle={Cosmo-Local Credit (CLC) White Paper v0.8},
  pdfauthor={William O. Ruddick and Mohamed Sohail}
}

\title{\foreignlanguage{english}{Cosmo-Local Credit (CLC)}\\شبكة لتوجيه الائتمان وتنسيق الالتزامات وتمويل اقتصاد عالمي ومحلي سليم}
\author{William O. Ruddick \\ Mohamed Sohail \\ Grassroots Economics Foundation \\ \texttt{info@grassecon.org}}
\date{White Paper v0.8 translation: 10 October 2026}

\begin{document}
\maketitle
\tableofcontents
\newpage
\sloppy
\RaggedLeft

\textbf{حول هذه الترجمة.} هذه ترجمة لـ White Paper v0.8. نُشر النص الإنجليزي المصدر في 30 سبتمبر 2026، ونُشرت هذه الترجمة في 10 أكتوبر 2026. الإنجليزية هي لغة النص المصدر. \href{https://docs.cosmolocal.credit/white-paper}{اقرأ المصدر الإنجليزي}.

\textbf{الإصدار:} 0.8

\textbf{التاريخ:} 30 سبتمبر 2026

\textbf{المؤلفان:} William O. Ruddick وMohamed Sohail

\textbf{للتواصل:} \href{mailto:info@grassecon.org}{info@grassecon.org}

\textbf{التنزيل:} \href{https://docs.cosmolocal.credit/white-paper/Cosmo-Local-Credit-CLC-White-Paper-v8-ar.pdf}{ملف PDF للورقة البيضاء CLC v0.8 بالعربية}

\textbf{حالة النشر:} إصدار نهائي للنشر العام. لا تمثل هذه الورقة نصيحة استثمارية ولا عرضًا لبيع ورقة مالية أو أداة مالية ولا طلبًا لشرائها. وقد تراجع الإصدارات اللاحقة التصاميم الموضحة هنا.

\section*{كيفية قراءة هذه الورقة}
\addcontentsline{toc}{section}{كيفية قراءة هذه الورقة}

تجمع هذه الورقة أربعة أنواع من المواد. وتنطبق الحالة في جميع المواضع حتى عندما لا يكررها كل قسم.

\begin{longtable}{P{0.30\linewidth} P{0.62\linewidth}}
\toprule
الحالة & المعنى \\
\midrule
\endhead
\textbf{حالي} & سلوك متحقق منه في CLC App أو في العقود العامة لـ Protocol v1.1.0. ويُعد \href{https://docs.cosmolocal.credit/ar/protocol/overview}{مرجع البروتوكول} المصدر الواقعي لهذه العقود. \\
\textbf{معتمد على التنفيذ} & قدرة لا تتوفر إلا عندما يتيحها تنفيذ أو صندوق أو مقدم خدمة أو ولاية قضائية أو جهة حوكمة أو سياسة منشورة بعينها. \\
\textbf{تاريخي} & دليل أو وظيفة من خدمة Sarafu Network السابقة. ولا يُعد ذلك تبنيًا حاليًا لـ CLC أو دليلًا على انتقال مستخدم أو أصل أو سجل أو التزام. \\
\textbf{مقترح} & تصميم أو نموذج اقتصادي أو آلية حوكمة أو بند في خارطة الطريق لا يُعرض باعتباره منشورًا حاليًا. \\
\bottomrule
\end{longtable}

يشمل Protocol v1.1.0 الحالي التنفيذ المباشر عبر \texttt{SwapPool}، والتنسيق والتقييم الاختياريين، وسقوف أرصدة الرموز في الصندوق، ورسوم الصندوق، ورسم بروتوكول إضافيًا، و\texttt{SwapRouter} لعرض الأسعار فقط. أما التنفيذ متعدد الخطوات وتوجيه HTLC أو الضمان، والتسوية الصافية، والتأمين المشترك، ورمز حوكمة CLC المقترح، وصندوق شبكة CLC المقترح، وآليات الائتمان الرسمية، وبرامج سيولة الشبكة، فهي مقترحة أو معتمدة على التنفيذ.

\textbf{الجمهور المقصود:} البناؤون والمشرفون، ومهندسو البروتوكول ومدققوه، وممولو السيولة والضمانات، ومصممو الحوكمة، والمراجعون القانونيون ومراجعو الامتثال.

\section*{نظرة عامة}
\addcontentsline{toc}{section}{نظرة عامة}

Cosmo-Local Credit عائلة منتجات مبنية حول \textbf{بروتوكول تجميع الالتزامات (CPP)}. يصف CPP كيف يمكن لصناديق الالتزامات ذات الحوكمة المستقلة أن تنسق الالتزامات القابلة للاسترداد، وتنشر طرق أسعار الصرف والحدود، وتحتفظ بالمخزون، وتتيح التبادل الخاضع للمساءلة.

\textbf{صندوق الالتزامات} ترتيب خاضع للحوكمة. أما \textbf{\texttt{SwapPool}} الحالي في Protocol v1.1.0 فهو مكوّن عقدي واحد: خزنة رموز ومحرك مبادلة مباشرة. ويمكن لتنفيذ ذلك العقد الاتصال بسجلات وعارضي أسعار وسقوف لأرصدة الرموز وسياسات رسوم وخدمات أخرى.

يعني «كوسمو-لوكال» مشاركة المعايير والبرمجيات والمعرفة المفتوحة عالميًا، مع إبقاء الإصدار والوفاء والحوكمة والمساءلة في الواقع محليةً. وتظل الجهة المصدرة مسؤولة عن قسيمتها، ويظل مشرف الصندوق مسؤولًا عن قواعد صندوقه وعن أي ضمان يتعهد به صراحةً. ولا يضمن CLC App أو GEF تلقائيًا قسيمة أو صندوقًا أو قيمة أو مسارًا أو وضع سيولة أو نتيجة.

\section*{التنفيذ الحالي}
\addcontentsline{toc}{section}{التنفيذ الحالي}

تدير GEF تطبيق CLC App على \href{https://cosmolocal.credit}{cosmolocal.credit}. ويوفر الوصول إلى الحساب والمحفظة وكتالوج السوق العام وواجهات مدعومة للرموز والقسائم والعروض وصناديق الالتزامات والتحويلات والمبادلات المباشرة عبر الصناديق. ويعتمد التوفر على الواجهة والعقود المنشورة والمخزون والحدود والرسوم المضبوطة وحالة الشبكة والأهلية ومقدمي الخدمات والولاية القضائية والشروط المنشورة للجهة المصدرة أو الصندوق.

لا يجعل تشغيل الواجهة GEF بمفرده جهة مصدرة أو مشرف صندوق أو أمين حفظ أو مُقرضًا أو مقترضًا أو وسيطًا أو ضامنًا أو جهة استرداد أو مستشارًا أو شركة تأمين أو طرفًا في الالتزامات بين المشاركين. وتحكم \href{https://docs.cosmolocal.credit/ar/governance/terms}{شروط الخدمة} استخدام التطبيق.

يفصل Protocol v1.1.0 بين الأدوار المسؤولة والتحكم التقني. فقد يكون مشرف الصندوق ومالك \texttt{SwapPool} ومدير الوكيل والمتحكم في المكونات التابعة ومشرف الكتالوج ومستلم الرسوم أطرافًا مختلفة. راجع \href{https://docs.cosmolocal.credit/ar/introduction/concepts}{المفاهيم والمصطلحات} للمصطلحات المشتركة.

\section*{الأصل التاريخي}
\addcontentsline{toc}{section}{الأصل التاريخي}

سبقت Sarafu Network التاريخية CLC App وقدمت أدلة عن العملات المجتمعية وتجميع الالتزامات. ولم يؤد انتقال الخدمة إلى Cosmo-Local Credit إلى نقل كل حساب أو محفظة أو رمز أو قسيمة أو صندوق أو رصيد أو تقرير أو مشارك أو التزام تاريخي.

تستخدم المؤشرات التاريخية في هذه الورقة مجموعة بيانات Sarafu المتسقة داخليًا. وهي ليست أرقامًا للاستخدام الحالي لـ CLC. راجع \href{https://docs.cosmolocal.credit/ar/white-paper/executive-summary}{الملخص التنفيذي} للمصدر وفترة القياس.

\section*{نموذج التصميم}
\addcontentsline{toc}{section}{نموذج التصميم}

يستخدم CPP أربع وظائف عامة:

\begin{itemize}[leftmargin=*]
\item \textbf{التنسيق:} تحديد القسائم أو الأصول الأخرى المقبولة.
\item \textbf{التقييم:} نشر الطريقة المستخدمة لإنتاج أسعار الصرف أو عروض الأسعار في الصندوق.
\item \textbf{وضع الحدود:} تطبيق سقوف أرصدة الرموز الحالية أو ضوابط مخاطر أخرى تُنفذ بصورة منفصلة.
\item \textbf{التبادل:} إدارة المخزون وتنفيذ مبادلات الصندوق وفق الرسوم والحدود المنشورة.
\end{itemize}
لا يجوز لتنفيذ أن يضيف خدمات التوجيه أو المراقبة أو دعم السيولة أو الضمانات أو الاحتياطيات أو التأمين أو الحوكمة إلا وفق سياسات منشورة بصورة منفصلة. ولا يثبت الإدراج أو التنسيق أو وجود حد أو احتياطي أو سجل على سلسلة الكتل بمفرده قدرة الجهة المصدرة أو وقوع الوفاء أو وجود تأييد أو ضمان أو تأمين.

\section*{القيم ومبادئ التصميم}
\addcontentsline{toc}{section}{القيم ومبادئ التصميم}

\begin{itemize}[leftmargin=*]
\item \textbf{رعاية الناس:} دعم رفاه الأفراد والجماعات.
\item \textbf{رعاية البيئة:} الحد من الضرر البيئي ودعم الممارسات التجديدية.
\item \textbf{العدالة:} تعزيز الوصول الآمن والمنصف إلى الموارد والمعرفة والرعاية.
\item \textbf{المعاملة بالمثل:} تقاسم المخاطر والتكاليف والمسؤوليات والفوائض وفق قواعد معلنة.
\item \textbf{عدم الهيمنة:} مقاومة تركز السيطرة على الناس والبيانات والتمويل والمعرفة والموارد المادية.
\item \textbf{المرونة:} مساعدة المجتمعات على الاستعداد للاضطرابات الاقتصادية والسياسية والمناخية والتكيف معها.
\item \textbf{المساءلة المحلية:} إبقاء الوفاء في الواقع والمسؤولية القانونية لدى الأطراف المحددة.
\item \textbf{قابلية الخروج:} السماح للمجتمعات والتنفيذات المتوافقة بمغادرة السجلات أو الخدمات المشتركة من دون حذف الأرصدة أو الالتزامات الصحيحة.
\end{itemize}
الطبقة الرقمية بنية تحتية مشتركة للذاكرة والتنسيق. وليست بديلًا عن العلاقات أو أداء الجهة المصدرة أو الحوكمة المحلية أو المساءلة القانونية.

\section*{التدفقات الحالية والمقترحة}
\addcontentsline{toc}{section}{التدفقات الحالية والمقترحة}

\subsection*{التدفق المباشر الحالي}
\addcontentsline{toc}{subsection}{التدفق المباشر الحالي}

\begin{enumerate}[leftmargin=*]
\item يراجع حامل القسيمة القسيمة وجهتها المصدرة وشروطها والعروض المرتبطة بها.
\item يقبل الصندوق الأصول المدعومة ويعلن إعداداته والسلطات المسؤولة عنه.
\item يحصل الحامل على عرض سعر ويمنح التفويض اللازم لمبادلة مباشرة عبر الصندوق.
\item ينفذ \texttt{SwapPool} تسوية المبادلة على السلسلة ويحسب رسم الصندوق ورسم البروتوكول.
\item إذا اختار الحامل لاحقًا \textbf{«استرداد»} في التطبيق، يجهز التطبيق تحويلًا إلى مالك الرمز بوصفه تقديمًا للقسيمة للاسترداد.
\item تفي الجهة المصدرة بصورة منفصلة بما وعدت به وتسجل الإبراء حتى لا يُعاد استخدام الوحدات التي تم الوفاء بها.
\end{enumerate}
الحصول على مخزون من صندوق هو مبادلة. ولا يكون استردادًا إلا إذا تولى الصندوق بصورة منفصلة دور الجهة المصدرة ووفى بشروط القسيمة المعمول بها.

\subsection*{التصميم المقترح الأوسع}
\addcontentsline{toc}{subsection}{التصميم المقترح الأوسع}

يستكشف التصميم المقترح توجيه التنفيذ والتسوية الصافية والخدمات المشتركة وبرامج سيولة الشبكة وسياسات التأمين وأصول الحوكمة. ويتطلب كل عنصر تنفيذًا وأطرافًا مسؤولة وقواعد معتمدة وتمويلًا وضوابط وإفصاحًا عامًا قبل تطبيقه.

قد يستخدم نموذج الرسوم المقترح حصةً للشبكة مقتطعة من رسوم الصندوق ورسومًا مستقلة للتوجيه أو الخدمات. ويختلف هذا الاقتراح عن Protocol v1.1.0، حيث يُضاف رسم بروتوكول اختياري إلى رسم الصندوق.

لا يحصل الداعمون أو مقدمو السيولة على حقوق تلقائية في حصص صندوق أو السداد أو السحب أو المكافآت أو الحوكمة من \texttt{SwapPool} الحالي. ويجب أن يفصح أي برنامج حالي أو مستقبلي بصورة منفصلة عن نوع التحويل والطرف المسؤول والتحكم وحقوق الاسترداد أو السحب والخسائر والمكافآت والرسوم وحقوق الحوكمة.

\section*{خريطة القراءة}
\addcontentsline{toc}{section}{خريطة القراءة}

\begin{itemize}[leftmargin=*]
\item المصطلحات الحالية ودورة حياة الإجراءات: \href{https://docs.cosmolocal.credit/ar/introduction/concepts}{المفاهيم والمصطلحات}
\item العقود المنشورة والمتحقق منها: \href{https://docs.cosmolocal.credit/ar/protocol/overview}{Protocol v1.1.0}
\item الأدلة التاريخية: \href{https://docs.cosmolocal.credit/ar/white-paper/executive-summary}{الملخص التنفيذي}
\item الاتحاد وقابلية التشغيل البيني: الفصول 5 و8 و11
\item الضمانات المقترحة وحدود التأمين: الفصلان 10 و11
\item نماذج السيولة والرسوم المقترحة: الفصول 7 و9 و12
\item إطار القياس المقترح: الفصل 14 والملاحق أ-ج
\item خارطة الطريق المقترحة: الفصل 15
\item الاعتبارات القانونية والامتثال: الفصل 17
\end{itemize}
تُحفظ \href{https://docs.cosmolocal.credit/ar/white-paper/archive}{الإصدارات السابقة من الورقة البيضاء} بوصفها منشورات تاريخية استُبدلت بوضوح فقط.

\section*{الملخص التنفيذي}
\addcontentsline{toc}{section}{الملخص التنفيذي}

المملكة \textbf{بروتوكول تجميع الالتزامات (CPP)} يصف كيفية إدارة صناديق الالتزامات المستقلة بإدارة الالتزامات القابلة للاسترداد، ونشر أساليب أسعار الصرف والحدود، والاحتفاظ بمخزونات، وتمكين التبادل المسؤول. لا يزال الجهات المصدرة مسؤولين عن التزاماتهم بالبطاقة. لا تزال مشرفو الصناديق مسؤولة عن قواعد المجموعة وأي ضمانات تتضمنها صراحة. لا يعد CLC App ولا GEF ضامنًا عالميًا.

توفر Protocol v1.1.0 كتلة بناء للتبادل المباشر الحالية: \texttt{GiftableToken}، \texttt{SwapPool}، السجلات الاختيارية و وحدات التقييم، حدود توازن رموز المجموعة، رسوم المجموعة، رسوم بروتوكول إضافية، ومبلغ اقتباس فقط \texttt{SwapRouter}. العقود الحالية لا تسجل الوفاء في العالم الحقيقي، أو تخلق أسهم مُقدمة السيولة تلقائيًا، أو تنفيذ طرق متعددة الأوقات، أو توفير احتياطيات أو ضمانات أو تأمينات عالمية.

التاريخية \textbf{Sarafu Network} استخدمت مفاهيم مشاركة الالتزامات عبر العملات المجتمعية ومجموعات التوفير ونظم المساعدات المتبادلة والتزامات الإنتاج. بعد ذلك انتقلت خدماتها إلى CLC App. لم تكن هذه تمثيلًا أن كل حساب تاريخي أو محفظة أو رمز أو قسيمة أو مجموعة أو رصيد أو تقرير أو مشارك أو التزام هاجر.

\subsection*{صورة تاريخية Sarafu Network --- نشاط Celo من 5 يوليو 2023 إلى 20 يوليو 2025}
\addcontentsline{toc}{subsection}{صورة تاريخية Sarafu Network --- نشاط Celo من 5 يوليو 2023 إلى 20 يوليو 2025}

\textbf{المصدر:} \href{https://dune.com/grassrootseconomics/sarafu-network}{لوحة التحليلات Dune}

\begin{itemize}[leftmargin=*]
\item 26367 مستخدم
\item 285,197 تبادل بين الزملاء
\item 188 النشطة الفريدة صناديق الالتزامات
\item 745 البطاقات النشطة الفريدة
\item \$320,692 حجم تبادل المجموعة
\item 899 تقارير نشرت (\href{https://sarafu.network/reports}{أرشيف التقارير التاريخية})
\end{itemize}
هذه الأرقام هي أدلة تاريخية، وليس أرقام استخدام CLC الحالية.

يبحث التصميم الأوسع المقترح لـ CLC كيف يمكن لشبكات متوافقة:

\begin{enumerate}[leftmargin=*]
\item اكتشاف وسائل الاقتباس عبر السباحات التي تحكمها بشكل مستقل
\item تنسيق خدمات الشبكة المسؤولة دون أن تتداخل المسؤولية المحلية.
\item دعم تفويضات السيولة أو الضمانات أو الاحتياطيات أو سياسات التأمين الممولة بشكل منفصل؛
\item قياس تبادلات المجموعة بشكل منفصل عن تقديم البطاقات والتوفير والإبراء؛ و
\item استخدام رسوم الشبكة المقترحة ورسوم الخدمة لتمويل الخدمات المشتركة المعتمدة.
\end{enumerate}
ويشمل الاقتراح \textbf{إشارة حكم CLC المقترحة} و a \textbf{تجمع شبكة CLC المقترح}. . ولا هي جزء من التطبيق الحالي أو Protocol v1.1.0.

بموجب النموذج المقترح، يمكن للمشاركين في السيولة والحوكمة أن يتصرفوا فقط من خلال البرامج المعتمدة بشكل منفصل مع الأهلية المنشورة، والمخاطر، والسيطرة، وحقوق النقل، وشروط الاسترداد، وقواعد الرسوم. إيداعات المجموعة العادية لا تخلق أي حصة تلقائية، سداد، سحب، مكافأة، أو الحق في الحكم.

الهدف المشترك هو:

زيادة التبادل المسؤول والوفاء بالتزامات العالم الحقيقي مع الحفاظ على الرعاية والعدالة والمسؤولية المحلية والمرونة.

مكافحة الاعتقال والخروج الموثوق به ما زالوا أهداف التصميم الأساسية. وتشمل المراقبة المقترحة الحوكمة المتأخرة في الوقت المحدد، وعروض الموافقة المتعددة، والوكالة الشفافة، وقواعد النزاع، ومراجعة الحوادث، وعملية الموثقة لتنفيذ المشاكل. هذه الرقابة ستقلل من مخاطر الحوكمة المختارة ولكن لا يمكن القضاء على الخسائر أو ضمان النتائج.


\section*{1. بروتوكول تجميع الالتزامات (CPP): البدائي الأساسي}
\addcontentsline{toc}{section}{1. بروتوكول تجميع الالتزامات (CPP): البدائي الأساسي}

\textbf{النموذج العقلي:} إن صندوق الالتزامات هو ترتيب يحكم على قبول البطاقات أو الأصول الأخرى، ونشر قواعد الصرف، وعقد المخزون، وتمكين التبادلات. يقدم أصحاب البطاقات لاحقاً إلى أصحابها من أجل تحقيقها في العالم الحقيقي. تبادل المجموعة والوفاء بالجهة المصدرة دورات حياة منفصلة.

تقوم CPP بتنسيق القيمة من خلال الالتزامات الموصوفة بوضوح. يتم مناقشة النموذج في \href{https://willruddick.substack.com/p/grassroots-economics-the-book-is}{الاقتصاد الجذري: التفكير والممارسة}.

\subsection*{1.1 ما هو الالتزام؟}
\addcontentsline{toc}{subsection}{1.1 ما هو الالتزام؟}

الالتزام هو وعد الطرف المحدد بالتسليم المستقبلي، على سبيل المثال الغذاء أو النقل أو العمل أو التخزين أو أي سلعة أو خدمة أو فائدة أو أداء قانوني آخر. (أ) \textbf{البطاقة} هو علامة أو سجل تمثل هذا الالتزام بموجب الشروط المنشورة.

عقد الرمز يسجل الآليات الرقمية. تشير شروط البطاقة إلى الجهة المصدرة، العرض، القدرة، المكان، التوقيت، القيود، التقديم، الوفاء، الشكاوى، وعملية الإبراء.

\subsection*{1.2 ما هو صندوق الالتزامات؟}
\addcontentsline{toc}{subsection}{1.2 ما هو صندوق الالتزامات؟}

صندوق الالتزامات هو الترتيب المحكم. يمكن أن يتم إدارتها من قبل فرد أو تعاونية أو مجموعة مجتمعية أو وكالة عامة أو اتحاد أو شركة متعددة الأطراف أو مشغل خدمات أو بنية أخرى مسؤولة.

وتشمل الأدوار والسلطات ذات الصلة:

\begin{itemize}[leftmargin=*]
\item \textbf{مشرف الصندوق:} ينشر وإدارة قواعد المجموعة وأي ضمان يتم فرضه صراحة.
\item \textbf{مالك الصندوق:} يحمل صلاحيات مالك \texttt{SwapPool} الحالية.
\item \textbf{مدير الوكيل:} يمكن تحسين تنفيذ مؤقت؛
\item \textbf{مراقبي الاعتماد:} تحكم السجلات المكوّنة أو القوائم، أو القيود، أو مكونات الرسوم؛
\item \textbf{الباحث عن الطريق أو المشغل:} يمكن اكتشاف الاقتباسات أو، في تنفيذ مستقبلي، تنفيذ طريق مصرح به بشكل منفصل. و
\item \textbf{الضامن:} يتحمل الالتزام المحدد فقط من خلال الشروط المنشورة الممولة.
\end{itemize}
مجموعات CPP وظائف المجموعة إلى أربعة مفاهيم:

\begin{itemize}[leftmargin=*]
\item \textbf{الوصف:} اعترض على رموز أو قسائم معتمدة.
\item \textbf{التقييم:} يعلن عن الطريقة المستخدمة لسعر الصرف أو الاقتباسات.
\item \textbf{القيود:} تنطبق القيود الحالية لميزانية الرمزية في المجموعة أو غيرها من الضوابط التي يتم تنفيذها بشكل منفصل.
\item \textbf{التبادل:} الحفاظ على المخزون، تنفيذ عمليات تبادل، حساب الرسوم، وإصدار سجلات المعاملات.
\end{itemize}
يقوم Protocol v1.1.0 بتنفيذ هذه الوظائف من خلال \texttt{SwapPool} والاعتمادات الاختيارية. "\texttt{Limiter}" حالياً يحد من رصيد الرمز في مجموعة؛ لا توفر الحدود المتداولة، أو لكل حساب، أو على مستوى الشبكة. يُحسب \texttt{SwapRouter} حالياً الأسعار المتعددة المجموعات؛ لا تنفذ عمليات تبادل.

يمكن أن يضيف التصميم الأوسع المقترح لـ CPP حدود التداول والتحكم في الحسابات والمركبات التنفيذية و HTLC أو مسارات الاحتفاظ بها، وشبكة اللحوم. هذه هي المكونات المقترحة، وليس وصف للمحدد الحالي أو جهاز التوجيه.

\subsection*{1.3 منطق التبادل المباشر الحالي}
\addcontentsline{toc}{subsection}{1.3 منطق التبادل المباشر الحالي}

التبادل المباشر الحالي للمجموعة:

\begin{enumerate}[leftmargin=*]
\item يتحقق من السجل الاختياري لعلامات الدخول والخروج؛
\item قياس المدخلات المستلمة؛
\item يحصل على اقتباس من المؤشر الموضح أو يطبق موازنة الوحدة الخام؛
\item يتحقق من رصيد رموز المجموعة الناتجة مقابل المحدد الاختياري؛
\item يحسب رسوم المجموعة وأي رسوم بروتوكول إضافية؛
\item يتحقق من المخزون المتاحة للخروج؛
\item يحول رسوم البروتوكول والمخرجات وحسابات رسوم المجموعة؛ و
\item إصدار أحداث تبادل.
\end{enumerate}
إن الاقتباس هو معايير المعاملة، وليس دليل على قدرة الجهة المصدرة أو قيمة الاسترداد أو القيمة العادلة أو قابلية التحويل النقدي أو ضمان.

\subsection*{1.4 استخدام أوسع}
\addcontentsline{toc}{subsection}{1.4 استخدام أوسع}

يمكن لـ صناديق الالتزامات دعم التبادل المجتمعي والإنتاج والمساعدة المتبادلة والبرامج العامة وغيرها من الهياكل المسؤولة. يمكن لمنتج ائتماني مُوثق بشكل منفصل استخدام قسيمة كضمان أو أداة سداد، لكنها تتطلب شروطًا إضافية وخاصة المعاملة. الإرسال العادي أو إيداع المجموعة أو تبادل المجموعة أو تقديم القسيمة للاسترداد أو الوفاء أو الإبراء ليس بطبيعة الحال قرض أو استرداد.

تستهدف CPP سجلات التداول المسؤولة و سجلات المعاملات التي يمكن مراجعتها ، وليس التداول المضاربي. سجلات السلسلة لا تزال لا تثبت تحقيق العالم الحقيقي أو التأثير الاجتماعي.


\section*{2. التحول المحاسبي: من الأصول إلى الثقة}
\addcontentsline{toc}{section}{2. التحول المحاسبي: من الأصول إلى الثقة}

التمويل التقليدي يبدأ \textbf{الأصول - الخصوم = الأسهم}. . إعادة تشكيل CPP هذا لاقتصاد الالتزام:

\begin{itemize}[leftmargin=*]
\item \textbf{قدرة القبول:} كم من الالتزامات التي يتعهد بها الجهة المصدرة لا يزال المجموعة أو الشبكة على استعداد لقبولها.
\item \textbf{الالتزامات المتبقية:} وعود لم يتم الوفاء بها من قبل الآخرين
\item \textbf{القدرة على الوفاء:} القدرة المثبتة للمصدر على الوفاء بهذه الوعود.
\end{itemize}
\subsection*{العلاقة بين الالتزام والقدرة:}
\addcontentsline{toc}{subsection}{العلاقة بين الالتزام والقدرة:}

قدرة القبول - الالتزامات المتبقية = قدرة القبول المتبقية

هذه الهوية المفاهيمية يمكن أن تدل على حدود المخاطر في المجموعة: الإصدار غير المحدود لا يعني قبول غير المحدود. ليست هوية الموازنة أو بيان بأن كل قسيمة هي قانونيا ديون أو قرض.

\textbf{مثال:} يقوم الناقل بإصدار 100 رحلة in في البطاقات. الصندوق يقبل ما لا يزيد عن 40 رحلة في وقت واحد. إذا كان 15 رحلة مقبولة لا تزال مفقودة، فإن المجموعة لديها القدرة على قبول ما يصل إلى 25 رحلة إضافية بموجب هذه السياسة. الحساب نفسه لا يخلق قرض أو ضمان.



\section*{3. نشاط الوفاء والإبراء والتبادل}
\addcontentsline{toc}{section}{3. نشاط الوفاء والإبراء والتبادل}

تغيرات تبادل المجموعة التي تعالج الأصول المحتملة. تنفيذ البطاقة يحدث عندما يؤدي الجهة المصدرة الالتزام المنشور. تسجل الإبراء الوحدات المكتملة بحيث لا يمكن عرضها مرة أخرى. هذه الأحداث لا يجب أن تتراجع إلى استخدام واحد للتسوية.

الإطار المقترح للقياس يستخدم سجلات منفصلة ل:

\begin{itemize}[leftmargin=*]
\item حجم تبادل المجموعة المكتملة
\item تقديمات فدية سارية؛
\item الوفاء المؤكد للمصدر
\item الإبراء؛
\item الالتزامات المؤهلة المعلقة؛ و
\item تم قياس مخزون الصندوق
\end{itemize}
السرعة المقترحة لإفراج الالتزامات يمكن مقارنة القيمة المكتملة خلال فترة مع متوسط قيمة الالتزامات المؤهلة المتبقية، ولكن فقط عندما تستخدم كلاهما نفس طريقة التقييم المعلنة. إمدادات العملات الرقمية ليست مصطلحًا كافًا: يمكن أن تشمل مخزون الجهة المصدرةات أو الوحدات المنتهية من الصلاحية أو غير الوصول إليها، والاختبارات، والتوازنات التي لا تمثل التزامًا مستمرًا للأطراف الثالثة.

يُمكن مقارنة قيمة تبادل المجموعة المكتملة مع مخزون المجموعة المقياس مع نسبة فعالية تبادل المجموعة المختلفة. إنه يصف نشاط التبادل وليس الوفاء للمصدر.

إن السيولة يمكن أن تحسن توافر المخزونات وتسهيل التبادل. لا تضمن أن الإصدارات ستنجح، أو أن المساهمين سوف يستردون الأصول، أو أن اقتباس المجموعة يمثل استرداد أو قيمة نقدية. جميع حقوق مقدم السيولة تتطلب شروط نشرت منفصلة.

انظر \href{https://docs.cosmolocal.credit/ar/white-paper/appendix-a-math-box}{الملحق (أ)} للتعريفات و \href{https://docs.cosmolocal.credit/ar/white-paper/appendix-c-kpi-definitions}{الملحق (ج)} لمواصفات القياس المقترحة.


\section*{4. ضمان مستقبلي قابل لإعادة الاستخدام}
\addcontentsline{toc}{section}{4. ضمان مستقبلي قابل لإعادة الاستخدام}

قد تقبل سهولة الائتمان الموثقة بشكل منفصل قسائم قابلة للتحويل كضمان أو كوسيلة سداد. إذا كان كذلك

\begin{itemize}[leftmargin=*]
\item إدخال المجموعة يمكن أن يجعل الضمان أكثر قابلية للاكتشاف
\item خيارات إضافية لتقديم القانون والتنفيذ يمكن أن تدعم السداد؛
\item ستبقى الحدود، المخزون، التقييم، السيولة، أداء الجهة المصدرة، ومخاطر الإعاقة؛ و
\item لن يتم ضمان أي طريق أو استيفاء أو استرداد أو قيمة أو استرداد.
\end{itemize}
\subsection*{4.1 حلقة ائتمان المنتجين}
\addcontentsline{toc}{subsection}{4.1 حلقة ائتمان المنتجين}

يمكن لخدمة CLC- متوافقة في المستقبل دعم ائتمان المنتجين فقط عندما تقوم الأطراف بتأسيس سهولة منفصلة وتقدم صراحة المعاملة على أنها مقدمة قابلة للسداد. شروط المعاملة ستحدد الدائن والمدين وتقول:

\begin{itemize}[leftmargin=*]
\item مبالغ مقدمة واسترداد؛
\item المواعيد المتأخرة، والفوائد، والرسوم، واستخدام الضمان المسموح به؛
\item التخصيص وأي تغيير في الدائن أو صاحب المطالبة؛
\item كيف يتم تقييم وتثبيت الوفاء بالبطاقة لحساب الإسترداد؛
\item الإسترداد الجزئي، والمبالغ المتبقية، وعدم الوفاء، والإبراء. و
\item وسائل العلاج المتاحة بموجب القانون المعمول به.
\end{itemize}
التدفق الموضح هو:

\begin{enumerate}[leftmargin=*]
\item إن المنتج أو مقدم الخدمات يصدر قسيمة تمثل الالتزام بالإنتاج المستقبلي، مثل عشرة رحلات تاكسي، 50 كيلوغرام من الذرة، أو عشرة ساعات عمل.
\item تعترف مشرف الصندوق بهذا البطاقة وتعلن عن طريقة سعر صرف المجموعة، والحدود، والرسوم، وقواعد المخزون، وأي حماية يفترضها صراحة.
\item يقدم القرض أو برنامج السيولة بشكل منفصل مقدمة بموجب شروط المعاملة المكتوبة ويقبل قسائم المنتج كضمان أو أداة سداد متفق عليها.
\item يمكن لأصحاب البطاقات تحويل البطاقات، وتبادلها عبر المجموعات المتوافقة، أو تقديمها للمصدر.
\item الوفاء المؤكد من قبل الجهة المصدرة يلبي الالتزام بالبطاقة. إنه يقلل من رصيد الائتمان المنفصل فقط إلى الحد والقيمة وبناء على الأدلة المحددة في شروط الائتمان.
\item تتميز السجلات بين استيفاء قسائم من محاسبة الائتمان وتحدد أي استرداد جزئي أو مبلغ متبقي أو تفويض أو عدم الوفاء أو الإبراء.
\end{enumerate}
هذه الهيكل يمكن أن يسمح باسترداد الموافقة في النوعية مع الحفاظ على سجلات منفصلة لالتزامات البطاقة والائتمان. لا تنشأ من تحويل عادي أو إيداع أو تبادل المجموعة أو تقديم القسيمة للاسترداد أو الوفاء أو الإبراء.

\textbf{مثال مثالي:} يحصل المنتج على مقدمة قدرها 1000 دولار بموجب شروط تشير إلى أن الوفاء المؤكد بطاقات الذرة المحددة يُسجل مقابل مبلغ السداد باستخدام طريقة تقييم محددة. لا يطبق المقرض ائتمان إلا بعد استلام الأدلة المطلوبة بموجب هذه الشروط. إذا كانت شروط الائتمان لا تجعل هذا الارتباط، فإن استحواذ أو تحويل أو تبادل أو تقديم أو استيفاء قسيمة الائتمان لا يؤثر على مقدمة.


\section*{5. من صناديق منفصلة إلى شبكة اتحادية}
\addcontentsline{toc}{section}{5. من صناديق منفصلة إلى شبكة اتحادية}

تدعم Protocol v1.1.0 الحالية التنفيذ المباشر من خلال \texttt{SwapPool} و توفر \texttt{SwapRouter} بإقتباس فقط. فإنه لا يقوم بتنفيذ طرق متعددة المكالمات، أو HTLCs، أو طرق الاحتفاظ بالأمانة، أو تشبيك المجموعات، أو التصفية عبر الشبكة.

يُقترح هذا الفصل كيف يمكن للمجمعات الحكومية المستقلة التنسيق دون التخلي عن قواعد القبول والقياس والحد والرسوم والمخزون والترخيص والحوكمة الخاصة بها.

\subsection*{5.1 تدابير منفصلة للتبادل والتنفيذ}
\addcontentsline{toc}{subsection}{5.1 تدابير منفصلة للتبادل والتنفيذ}

الاتحاد قد يحسن الوصول إلى المخزون، لكنه لن يدمج البطاقات وتبادل دورات الحياة. أي تنفيذ يقيس هذه الأحداث بشكل منفصل:

\begin{enumerate}[leftmargin=*]
\item يتم الإشارة إلى طريق؛
\item يتم تنفيذ تبادلات مجموعة واحدة أو أكثر وتسوية في السلسلة.
\item يقدم حاملاً وحدات البطاقة للمصدر؛
\item يقوم الجهة المصدرة بالوفاء بالالتزام؛ و
\item يتم تفريغ الوحدات المكتملة.
\end{enumerate}
المزيد من الطرق المشار إليها أو التي تم تنفيذها لا تثبت تحقيقًا أكبر. وتشير التقارير إلى المجموعة، والفترة، والأصول، وأسلوب التقييم والخاتم الزمني، والاستبعادات، والتصحيحات، والأدلة خارج السلسلة المطلوبة في الملحق C.

\textbf{الطريق الموضح:} المدرسة لديها بطاقات الذرة ولكنها تحتاج بطاقات النقل خدمة الطرق تحدد المخزونات المتوافقة. سوف تنجح تنفيذها فقط إذا بقيت كل هوب مصرح بها بشكل منفصل ضمن حدود الاقتباسات والحدود والرسوم والمخزونات والسياسة. لن تثبت المبادلات المبادلة الناتجة أن أي من الجهات المصدرة قد أوفى لاحقاً بالتزاماته بالبطاقة.

\subsection*{5.2 الخدمات المقترحة للتوجيه وإعادة التوازن}
\addcontentsline{toc}{subsection}{5.2 الخدمات المقترحة للتوجيه وإعادة التوازن}

خدمة طريق مستقبلية يمكن أن تدعم نشاطين منفصلين.

\textbf{الإعدام الذي بدأه المشاركون} بالنظر إلى أصول الدخول والخروج، ومبلغ، وقيود المستخدمين، يمكن للخدمة تحديد مسار وإعداد التنفيذ. كل قفز سيكون لديه مسؤولية خاصة، الاقتباس، الموافقة، الرسوم، الحدود، المخزون، والإيصال. المجموعات الذرية، و HTLCs، و الاحتفاظ هي خيارات التنفيذ المستقبلية المحتملة، وليس سلوك البروتوكول الحالي.

\textbf{إعادة التوازن} يمكن لـ مشرفو الصناديق نشر أهداف المخزون، والعقود المسموح بها، وفرق الأصول، وحدود الانحراف من الاقتباسات، والحدات لكل فترة. يمكن للخدمة المسؤولة البحث عن دورات أو سلاسل متوافقة وتنفيذ فقط الأهداف المسموح بها.

إعادة التوازن سيكون اختياراً يمكن للمجموعة أن تسمح بطرق للمشاركين مع رفض إعادة التوازن الخارجي، أو يمكن أن تسمح فقط بأصول ومقارنات ومبالغ مختارة. كل صعود يتم تنفيذه سيؤدي إلى إيصال، ويتم الكشف عن أي رسوم خدمة بشكل منفصل عن رسوم المجموعة والبروتوكول.

\subsubsection*{5.2.1 الاتحاد والتفاعلية}

يمكن للمنشآت المستقلة تشغيل سجلاتها الخاصة، واجهاتها، وخدمات الطرق، وملفات السياسة في حين تختار بيانات متوافقة ومعايير الاستلام. لا يزال التنفيذ عبر الملفات الشخصية يعتمد على النشر.

ملف تعريف متوافق سيكون:

\begin{itemize}[leftmargin=*]
\item تحديد جذور سجلها ومشغلي الخدمات ومراقبينها والشروط المعمول بها؛
\item الإفصاح عن الأطراف المقابلة المسموح بها أو الممنوعة، والأصول، والمعدلات، والطرق؛
\item يطبق كل مجموعة مشاركة الترخيصات والقيود والرسوم وتقييدات المخزون؛
\item الحفاظ على الأدلة من الاقتباس إلى الاستلام في كل مكان؛ و
\item السماح لمجموعات وظيفية أخرى بمغادرة أو اختيار سجل آخر دون حذف الرصيد أو التزامات الجهة المصدرة.
\end{itemize}
يمكن للموافقة زيادة مسارات التبادل المتاحة وتقليل الاعتماد على سجل أو مشغل واحد. لا تجعل الشبكة، CLC App، GEF، أو مجموعة أخرى مسؤولة عن الوفاء بالجهة المصدرة.

\subsection*{5.3 النموذج المقترح للشبكة والرسوم الخدمية}
\addcontentsline{toc}{subsection}{5.3 النموذج المقترح للشبكة والرسوم الخدمية}

تكلفة Protocol v1.1.0 الحالية رسوم المجموعة، وعندما يتم تشكيلها، رسوم بروتوكول إضافية على تبادل المجموعة المباشر. هذه الرسوم الحالية لا تزال متميزة.

يمكن لبرنامج مستقبلي أن يتلقى بشكل منفصل:

\begin{enumerate}[leftmargin=*]
\item (أ) \textbf{الشبكة}تعريفها باعتبارها حصة معينة من رسوم المجموعة التي تم جمعها من المجموعات المشاركة . و
\item (أ) \textbf{رسوم التوجيه أو خدمة}، تتقاضى مقابل خدمة مستقبلية محددة
\end{enumerate}
لا تعتبر الرقصة المقترحة للشبكة نسبة مئوية إضافية تطبق على مبلغ التبادل الكامل بعد احتساب رسوم المجموعة بالفعل. بالنسبة لمجموعة \texttt{p}:

\[
\tau_p = f_p \cdot r_p
\]

حيث \texttt{f\_p} هو معدل رسوم المجموعة و \texttt{r\_p} هو الجزء المقترح من هذه رسوم المجموعة المخصصة لبرنامج الشبكة.

لفترة قياسية:

\begin{itemize}[leftmargin=*]
\item \textbf{رسوم المجموعة الإجمالية} هو مبلغ الرسوم الفعلية التي تم جمعها في كل مجموعة.
\item \textbf{الإيصالات المتعلقة بالشبكة} هي الأسهم المعلنة من هذه الرسوم المستحقة؛
\item \textbf{إيصالات رسوم الخدمة} يتم تحصيل رسوم التوجيه أو خدمة بشكل منفصل. و
\item \textbf{إيصالات رسوم البرنامج} الإيصالات المتساوية للشبكة بالإضافة إلى إيصالات رسوم الخدمة.
\end{itemize}
لا توجد فئة يتم احتسابها مرتين لا تتضمن رسوم البروتوكول الحالية ما لم يتم إعادة توجيه السياسة المعتمدة بشكل منفصل بشكل قانوني إيرادات رسوم البروتوكول الفعلية إلى البرنامج المستقبلي.

بالنسبة للمقارنة الإجمالية، دع:

\begin{itemize}[leftmargin=*]
\item يكون \texttt{Q\_swap} قيمة المقاييس المجموعة التي تم تنفيذها للفئة والمدة المحددة ؛ و
\item \texttt{$\tau$} هو المعدل الفعلي لضربة الشبكة المقترحة ومكافآت الخدمة المحددة بشكل منفصل على قيمة المقايضة المنفذة.
\end{itemize}
ثم:

\[
F \approx \tau \cdot Q_{\text{\latinfont swap}}
\]

هذه تقدير تحليلي وليس وعد بالدخل كل إدخال يتطلب مجموعة محددة، فترة، ووحدة، وخط زمني للتقييم، واستبعاد، وسياسة التصحيح.

\subsubsection*{5.3.1 الإيصالات القابلة للحصول على النقدية و الإيصالات الطبيعية}

قد تصل الرسوم إلى الأصول المزدوجة المؤهلة للحصول على النقود أو إلى البطاقات وغيرها من الأصول الطبيعية. لا يمكن للإيصالات الطبيعية دفع نفقات نقدية أو مطالبات التغطية تلقائيًا. أي تبادل أو تحويل يتطلب السلطة والكتابة المتاحة ومواقع الكشف عنها والحدود والتنفيذ الفعلي.

دعونا \texttt{$\chi$} تكون النسبة المحققة من إيرادات الرسوم التي تكون قابلة للحصول على النقدية بعد قيود السياسة والتحويلات الفاشلة والانزلاق. الإيصالات القابلة للاستخدام النقدي هي:

\[
F_{\text{\latinfont cash}} \approx \chi \cdot F
\]

تحليل الميزانية وتحليل التوازن يستخدم \texttt{F\_cash} المحقق ، وليس الرسوم الإجمالية المدرجة أو القيمة الاسمية للمخزون الطبيعي. في المستقبل، سيتم الإبلاغ عن رسوم المجموعة الإجمالية، وإيرادات الشبكة، وإيرادات الخدمات، وتكوين الأصول، ونتائج التحويل، وإيرادات القيمة النقدية التي يمكن استخدامها بشكل منفصل.

\subsection*{5.4 برامج السيولة المقترحة}
\addcontentsline{toc}{subsection}{5.4 برامج السيولة المقترحة}

يمكن لبرنامج السيولة المستقبلي الموثق بشكل منفصل تخصيص الأصول إلى المجموعات المعينة أو خدمات التوجيه. العقود الحالية \texttt{SwapPool} لا تضع أسهم الصندوق أو تخلق تلقائيًا حقوق السداد أو السحب أو المكافأة أو الحوكمة أو الربح.

أي برنامج ينشر:

\begin{itemize}[leftmargin=*]
\item الكيان المسؤول والمشاركة مشرفو الصناديق؛
\item الأصول المساهمة وما إذا كان التحويل قابلاً للرد أو قابلاً للسحب أو تبرع به أو تمنحه؛
\item ترتيبات الوصاية والتحكم الفني؛
\item الاستخدامات المسموح بها، والحدود، والقفلات، وبوابات السحب، وتخصيص الخسائر؛
\item تأهل الرسوم أو الحوافز وما إذا كان المبلغ قد يكون صفرًا؛
\item الإبلاغات والنزاعات والشكاوى والإجراءات العلاجية؛ و
\item الهجرة، وإنهاء، ومعاملة الأصول والالتزامات المتبقية.
\end{itemize}
وتشمل المخاطر المادية المخزونات التي يصعب تبادلها أو الوفاء بها، وانخفاض مؤهلية النقد، وعدم الوفاء بالجهة المصدرة، وفشل العقود أو مزودها، وتغييرات الحوكمة، والقيود على الخروج. قد تقلل الحدود والاحتياطيات والإيصالات وألواح التحكم أو تكشف عن بعض المخاطر. لا يزيلون الخسارة

بالنسبة لمقاييس تحليلية سابقة، دعونا:

\begin{itemize}[leftmargin=*]
\item يكون \texttt{$\phi$} الجزء الحقيقي من إيرادات رسوم البرنامج المخصصة بموجب شروط البرنامج المعتمدة؛ و
\item \texttt{K} هو القيمة المقياسة للأصول التي تغطي البرنامج بموجب طريقة واحدة محددة.
\end{itemize}
ثم:

\[
\text{\latinfont FeeFlow}_{LP} \approx \frac{\phi \cdot F}{K} = \frac{\phi \cdot \tau \cdot Q_{\text{\latinfont swap}}}{K}
\]

هذه المقياسة تصف تدفق الرسوم المحقق لكل أصول البرنامج المقياسة. إنها ليست APY، أو توقعات، أو أرباح، أو عائدات مضمونة. وتحتفظ التقارير بحجم المقايضة، وتقديم القسيمة للاسترداد، وتوفير الجهة المصدرة، والإبراء، ومدة الاحتفاظ، والخسائر، والسحب، وإيصالات الرسوم بشكل منفصل.


\section*{6. السيولة والحوكمة المقترحة على مستوى الشبكة}
\addcontentsline{toc}{section}{6. السيولة والحوكمة المقترحة على مستوى الشبكة}

التجربة من Sarafu Network التاريخية دفعت البحوث إلى سؤالين أوسع من التصميم:

\begin{enumerate}[leftmargin=*]
\item كيف يمكن للمجموعات الحكومية المستقلة تنسيق خدمات السيولة والتوجه ؛ و
\item كيف يمكن للخدمات المشتركة أن تبقى مسؤولة مع نمو المشاركة.
\end{enumerate}
يقدم أحد التصاميم المقترحة CLC \textbf{تجمع شبكة CLC المقترح} كمنظمة للتصفية على مستوى الشبكة وتنسيق الميزانية و \textbf{إشارة حكم CLC المقترحة}. . لا يوجد أي منها في التطبيق الحالي أو Protocol v1.1.0. يمكن أن تستخدم النشرات الأخرى التعاونيات، والوكالات العامة، والاتحادات، والشركات المتعددة، ومشغلي الخدمات، أو غيرها من الهياكل المسؤولة دون اعتماد أي من الاقتراحات.

\subsection*{6.1 التحكمات الجانبية والخيارات}
\addcontentsline{toc}{subsection}{6.1 التحكمات الجانبية والخيارات}

يمكن لـ Protocol v1.1.0 الحالية تطبيق حدود ميزان رموز الصندوق والتحقق من المخزون. هذه الضوابط تقيد إجراءات عقدية مختارة؛ فإنها لا تمنع عدم أداء الجهة المصدرة، وفقدان القيمة، وفقدان المفتاح، وفشل العقد، أو أي خسارة أخرى.

في المستقبل، قد يضيف التنفيذ قطع الدوائر، وقوانين التوقيت، احتياطيات، ضمانات، أو التأمين. أي حماية تحتاج إلى طرف مسؤول محدد، والحدث المغطى، والتمويل، والحد الأقصى، والإقصاءات، والمدة، والأدلة، وعملية المطالبات، والشروط المعمول بها.

يجب أن تبقي الحوكمة المقترحة صلاحيات التسجيل والخدمات ضيقة ومراجعة. الإجراءات العادية يمكن أن تستخدم إشعار، تأخير، نشر السبب، الاستئناف، وتصحيح العمليات. وينبغي أن توفر الإجراءات الطارئة سجلات الحوادث وتنتهي صلاحيتها أو تتم مراجعتها بموجب السياسة المعتمدة.

اكتشاف مسار متوافق يمكن أن يجعل المزيد من التبادلات ممكنة عبر الصندوق. إن تنفيذ المجموعات المتعددة وتسويقها يتطلب برنامجًا إضافيًا ، والإذن ، والحدود ، ومعالجة الفشل ، والحكم. لا يثبت المزيد من عمليات التبادل المشهورة أو المنفذة نفسها استيفاء البطاقة أو التأثير الاجتماعي.


\section*{7. اقترح CLC أصول الإدارة والحوكمة}
\addcontentsline{toc}{section}{7. اقترح CLC أصول الإدارة والحوكمة}

يقدم هذا الفصل تصميمًا اختياريًا للحوكمة المستقبلية وتنسيق السيولة. رمز الحوكمة CLC المقترح، stCLC، sCLC، خزانة الحوكمة، صندوق التأمين، آبشار، مقياسات، ولايات، نافذة الائتمان الرسمي، والبرامج السوق الخارجية ليست جزءا من CLC App الحالي أو Protocol v1.1.0. كل منهما يتطلب تنفيذًا ومشغلًا مسؤولًا، وسياسة وشروطًا منشورة، مراجعة أمنية، والموافقة القانونية المطبقة.

CPP لا تتطلب هذا النموذج الرمزي. يمكن أن تستخدم النشرات المتوافقة بدلاً من ذلك التعاونيات والوكالات العامة والاتحادات والشركات المتعددة والخدمات أو غيرها من الهياكل المسؤولة.

\subsection*{7.1 الغرض}
\addcontentsline{toc}{subsection}{7.1 الغرض}

إذا تم اعتمادها، فإن طبقة الحوكمة المقترحة يمكن أن:

\begin{itemize}[leftmargin=*]
\item تنسيق السجلات المشتركة وخدمات الطرق؛
\item تخصيص تفويضات السيولة المعتمدة عبر المجموعات التي تحكمها بشكل مستقل ؛
\item إدارة برنامج تغطية يتم تمويله اختياريًا ومحدد صراحة؛
\item الحفاظ على البنية التحتية والملاحظة المشتركة.
\item الاعتراف بالمساهمين بموجب القواعد المنشورة للصلاحية ومكافحة الألعاب؛ و
\item الحفاظ على المراجعة والإستئناف والشفافية والخروج الموثوق به مع نمو المشاركة.
\end{itemize}
\subsection*{7.2 النموذج المقترح لحوكمة الأصول}
\addcontentsline{toc}{subsection}{7.2 النموذج المقترح لحوكمة الأصول}

يحتوي التصميم على ثلاثة أصول مختلفة المقترحة:

\begin{enumerate}[leftmargin=*]
\item \textbf{تم اقتراح رمز الحكم CLC:} أصول حوكمة أساسية قابلة للتحويل بموجب قواعد الإصدار والاحتفاظ والتصويت والامتثال المعتمدة.
\item \textbf{stCLC:} إيصال صيانة التصويت غير قابلة للتحويل يتم صياغته عندما يتم قفل رمز الحوكمة CLC المقترح. وسوف تمثل سلطة حكم محدودة بالوقت ويمكن أن تسمح بالتفويض المكتشف.
\item \textbf{sCLC:} تصريح أو أصول تحفيزية تم إصدارها في إطار السياسة. يمكن أن تدعم الوصول إلى الائتمانات الرسومية المحدودة أو الحوافز المعتمدة للسيولة والتوجيه ويمكن أن تنتهي صلاحيتها أو يتم حرقها في نهاية الفترة.
\end{enumerate}
ولا stCLC ولا sCLC ستكون رأس المال، أو أرباح، أو مطالبة بقية، أو وعد بالعائد، أو قسيمة اجتماعية. السياسة المعتمدة يمكن أن تحدد إصدار sCLC والوصول إلى الرسوم الائتمانية إلى الصفر في أي عصر.

سيتم نشر قوائم الحكم، والمدة الحد الأدنى، والتخفيفات، والتفويضات، والحد الأقصى للتركيز، والحد الأدنى للعمل الحرج قبل الاستخدام. لا يمكن أن يصوت إشارات الحكم المقترحة CLC المفتوحة بموجب هذا التصميم.

\subsubsection*{7.2.1 إمدادات وتخصيصات مثالية}

الإجمالي التوضيحي المقترح CLC إمدادات رموز الحكم هو \textbf{500,000,000}. . هذه هي معايير التصميم، وليس إصدارًا أو عرضًا أو تقييمًا أو وعدًا بالسيولة.

تخصيص مثالي هو:

\begin{itemize}[leftmargin=*]
\item \textbf{15\% --- Grassroots Economics Foundation:} المشاركة بشكل دائم ، غير قابلة للتحويل ، ومرتبطة بـ GEF multisig المحدد بموجب قواعد التوقيع والدورة المنشورة.
\item \textbf{15\%  الفريق الأساسي والشركاء الأوائل:} المشاركة خلال فترة محددة للسياسة والفترة السطرية للاستحقاق لمدة 24 شهرا، مع الكشف مسبقا عن قواعد التصويت والتحويل.
\item \textbf{30\%  المنح الخاصة:} الاعتراف للمساهمين المبكرين الذين يقدمون سيولة الشبكة المعتمدة، مع الاستحقاق المحدد للسياسة لمدة 24 شهرا.
\item \textbf{40٪ --- احتياطي السيولة العامة:} يتم الاحتفاظ بها في صندوق إدارة مغلق في الوقت لبرامج السوق الخارجية والوصول الاختياري.
\end{itemize}
في إطار الضوابط الوضوحية للسيولة العامة:

\begin{itemize}[leftmargin=*]
\item لن يكون أكثر من 10\% من إجمالي العرض نشطًا في جميع الأماكن الخارجية في وقت واحد.
\item كل عملية نشر تنتهي بعد 90 يوما ما لم يتم تجديدها من خلال العملية المعتمدة.
\item سيظل موقف السيولة وأي رموز موقف تحت سيطرة الحوكمة المعلنة. و
\item إشارات الحوكمة المقترحة CLC المستخدمة في مواقف السيولة لن تصوت إلا إذا تم حبسها بموجب نفس القواعد مثل إشارات التصويت الأخرى.
\end{itemize}
لا يمكن أن يبدأ إطلاق مستقبلي بإفصاح متعددة الألواح والانتقال إلا بعد موافقة منفصلة على أساسيات التصلب، مثل عمليات التدقيق المستقلة والمراقبة ومكتبات تشغيل الحوادث والإجراءات المختبرة للاستراحة والشوك. سيتم نشر كل معايير و عملية التغيير المعتمدة بشكل منفصل.

\subsubsection*{7.2.2 المراحل المتاحة والمنح}

يمكن لبرنامج المنح المستقبلي استخدام نوافذ المساهمة المرحلية والقيمة المرجعية المنشورة للدخل والميزانية. هذا الإشارة لن تكون وعد بسعر السوق أو التقدير أو السداد أو الخروج.

تنص شروط المساهمة على ما إذا كان الأصول قد تم التبرع بها، أو تمنحها، أو يمكن سحبها، أو يمكن تسديدها. الذي يسيطر عليهم استخدامها المسموح به. القفلات؛ تخصيص الخسائر الإبلاغ و شروط الخروج يمكن أن تواجه المساهمات الكبيرة حدود التصويت أو انخفاض وزن الحوكمة.

أي سيولة في السوق الخارجية تتطلب قرارًا منفصلًا يحدد المواقع والعاملين المسؤولين والأصول والحد الأقصى والوقت والصراعات ومراقبة المخاطر والتقديم للإبلاغ والإنهاء. السياسة لن تستهدف أو تضمن سعرًا.

\subsubsection*{7.2.3 البرنامج المقترح لتربية التأثير}

يمكن لبرنامج مستقبلي تخصيص جزء من صندوق الحوكمة للمساهمين الذين يضعون الأصول في صناديق الالتزامات المحدد ويحسنون بشكل قياسي الوصول إلى البورصة وتلبية الجهة المصدرة المؤكد.

سياسة التأهل المعتمدة:

\begin{enumerate}[leftmargin=*]
\item تحديد المجموعات المعتمدة، والأصول، الحد الأدنى من المدة، وشروط الإغلاق؛
\item تحديد المساهمة الإنتاجية باستخدام الأدلة المتعلقة بالمقايضة التي تم تنفيذها وإثباتها بشكل منفصل وتوفيرها وإفراجها؛
\item قياس المساهمة الهامشية بدلاً من القيمة الإجمالية المحددة وحدها
\item يستبعد أنشطة الغسيل الذاتية والدورية؛
\item تطبق أقصى حدود لكل كيان وتقليل العوائد؛ و
\item توفير نافذة للملاحظة والنزاع والتصحيح قبل التخصيص النهائي.
\end{enumerate}
ستظل أي منحة تخضع للشروط المنشورة للبرنامج، والحصول على المال، والقدرة على التأهل، والامتثال، وقواعد الخسارة.

\subsection*{7.3 اقتراحات التصويت والحوافز والوصول إلى الرسوم}
\addcontentsline{toc}{subsection}{7.3 اقتراحات التصويت والحوافز والوصول إلى الرسوم}

بموجب هذا التصميم، يمكن لأصحاب stCLC التصويت على المجموعات المؤهلة، أو ولايات خدمة الطرق، أو ميزانيات مشتركة، أو المقترحات الأخرى المعتمدة. يمكن أن ينقل نظام القياس المنتخب الأصوات إلى حوافز sCLC محدودة لمشغلي السيولة المؤهلة أو التوجيه.

سياسة القياس ستبقي المبادلات المبادلة التي تم تنفيذها والوفاء بالجهة المصدرة منفصلين. إنه لن يعطي مكافأة للطرق المذكورة ولكن غير المنفذة، والعروض غير المنفذة، أو التداول الذاتي، أو القيمة الإجمالية مقفلة دون دليل على النشاط المفيد.

يمكن أن ينشر عملية الحوكمة المعتمدة أيضا ميزانية إئتمان رسوم العصر. أصحاب stCLC المؤهلين قد يحصلون على sCLC الموافقة على التبادلات المحدودة، المحدودة في الوقت من المجموعات المحددة التي تحمل الرسوم إلى الأصول أو البرامج المسجلة. هذا سيكون الوصول المحدد للمخزون المتوفر، وليس الدخل السلبي أو الملكية في مجموعة الاحتفاظ بالرسوم.

\subsubsection*{7.3.1 إنفاذ رسوم الخدمة واختيار الملف الشخصي}

قد يتطلب ملف التسجيل المستقبلي أن تستخدم المجمعات المشاركة أو خدمات الممرات مُعدلًا أو مضربًا متوافقًا للرسوم. قد يرفض الملف الشخصي اكتشاف أو توجيه أو دعم البرنامج عندما لا يثبت الإيصالات المطلوبة جمع الرسوم.

أسماء مثل: \textbf{PoolFactory} أو \textbf{(فيه هوك)} وصف وحدات مستقبلية محتملة؛ هذه ليست عقود Protocol v1.1.0. أي تنفيذ سيكشف عن الرمز والعناوين والمراقبين والمستلمين والأسعار والمعاملات المؤهلة ومعالجة الفشل وحالة المراجعة.

التنفيذ سيكون محدداً يمكن لمجموعة مستبعدة من صفحة واحدة الاستمرار في العمل على المستوى المحلي أو الانضمام إلى سجل آخر متوافق إذا ظلت عقودها واعتماداتها تعمل. يجب على العملاء الكشف عن الملفات الشخصية المتاحة وإظهار الرسوم المطبقة قبل الحصول على الطلب.

\subsubsection*{7.3.2 الميزانية sCLC والتحكم الخارجي الاختياري في العبور}

بعد اعتماد التمويل ميزانيات ذات الأولويات العالية، يمكن في عملية مستقبلية نشر ميزانية إئتمان رسوم العصر \texttt{F\_epoch}، بما في ذلك الصفر. قد تحدد السياسة حدًا للمشاركين متناسبًا مع سلطة التصويت stCLC:

\[
\text{\latinfont limit}_{\text{\latinfont user,epoch}} = F_{\text{\latinfont epoch}} \times \frac{\text{\latinfont stCLC}_{\text{\latinfont user}}}{\text{\latinfont stCLC}_{\text{\latinfont total}}}
\]

ستظل كل نافذة خاضعة لقوائم الإجازة، والحدات، والمخزونات، والإجازة، وقواعد الحوادث، والقانون المعمول به.

سياسة منفصلة يمكن أن تسمح باستحواذ محدد معدل الوقت على رمز الحوكمة المقترح CLC فقط لتقليل سطح الهجوم على الحوكمة الخارجية المقياس. سوف يتطلب ذلك:

\begin{itemize}[leftmargin=*]
\item محفز المنشأ على أساس العجلة الخارجية على مدى فترة محددة.
\item حصة أقصى من الرسوم المحققة التي تكون قابلة للتأهيل وحجم الموقع؛
\item التنفيذ الموزن في الوقت دون إعلان عن التوقيت الدقيق؛
\item توقف تلقائي عند عدم استيفاء التغطية المعتمدة أو العدويات التشغيلية؛
\item التقاعد أو وضع الرموز المالية المكتسبة في حساب غير التصويت المعلن؛ و
\item بيان صريح بأن البرنامج لا يستهدف أو يضمن سعرًا ولا يؤثر على الوفاء بالجهة المصدرة.
\end{itemize}
السياسة يمكن أن تبقى معطلة إلى أجل غير مسمى

\subsubsection*{7.3.3 مجموعة المحافظيات والشهادات}

(أ) \textbf{مجموعة محفظة} سيكون صندوق الالتزامات العادي الذي يتم تنظيمه حول مجال المهمة أو الخدمة. يمكن أن يكون مشرف الصندوق فردًا أو تعاونًا أو مجموعة مجتمعية أو وكالة عامة أو اتحادًا أو مجموعة متعددة الأطراف أو مشغل خدمة أو بنية أخرى مسؤولة.

يمكن الحصول على الدعم من خلال:

\begin{enumerate}[leftmargin=*]
\item المساهمات المباشرة بموجب الشروط المنشورة للمجموعة.
\item تفويضات السيولة المحدودة في الوقت المحدد التي تمت الموافقة عليها من خلال عملية الحوكمة المعتمدة. أو
\item sCLC-الوصول الموجّه إلى ميزانية محدودة بعد سقوط المياه عندما يتمّ تمكين هذه الآلية.
\end{enumerate}
ستظل مجموعات المحافظين تحت إدارة مستقلة ويمكن استخدام سجلات مشتركة أو مستقلة.

يمكن أن تؤكد الشهادات على القبول أو معاملة المخاطر ولكن لن تخلق الحق في الرسوم أو الأرباح أو الأصول المتبقية. قد تستخدم النشر:

\begin{itemize}[leftmargin=*]
\item \textbf{شهادات مرتبطة بالتسجيل} من المؤكد المحدد الذي يطبق طريقة منشورة؛ أو
\item \textbf{بطاقات خدمة التحقق} تمثل التزامًا قابلًا للاسترداد لإجراء مراجعة أو تقييم مذكورة.
\end{itemize}
سوف يكشف المجموعة عن كيفية تأثير الشهادة على القبول، وأساليب سعر الصرف، والحدود، أو التأهل وكيف يتم التعامل مع الأخطاء، والإنتهاء، والنزاعات، والتقدم.

\subsection*{7.4 الاقتراحات والميزانيات}
\addcontentsline{toc}{subsection}{7.4 الاقتراحات والميزانيات}

سيقوم المشروع بتخصيص الإيرادات المؤهلة المحققة فقط في إطار عملية الحوكمة المعتمدة. سوف تمييز:

\begin{itemize}[leftmargin=*]
\item \textbf{الأصول القابلة للحصول على النقد (\texttt{E\_cash}):} الأصول المعدلة المدرجة التي يمكن استخدامها أو تحويلها بموجب سياسة التزامات النقدية. و
\item \textbf{الأصول الطبيعية (\texttt{E\_kind}):} البطاقات أو الأصول الأخرى التي قد تدعم البورصة في الشبكة ولكن لا تعتبر ضمن التغطية النقدية أو التزامات التشغيل ما لم يتم تحويلها فعليا.
\end{itemize}
سياسة التحويل ستحدد المشغلين المعتمدين، والأصول، والمواقع، ومصادر الأسعار، ونوافذ الزمن، وحدود الانزلاق، والحد الأقصى، والصراعات، والسجلات، وإجراءات الحوادث. إن تحويل الخزانة لا يحدد التزام البطاقات للمصدر أو طريقة سعر الصرف للمجموعة.

ترتيب الأولوية الموضح هو:

\begin{enumerate}[leftmargin=*]
\item \textbf{هدف التغطية الممولة:} بناء أي احتياطي تم تبنيه أو صندوق تأمين مقترح على هدفها المعلن عن الأحداث المغطاة المحددة.
\item \textbf{العمليات الأساسية:} تمويل ميزانية محددة ومعلنة للعمل القانوني والتعليم والاتصالات والبنية التحتية والتحققات والملاحظة.
\item \textbf{ولايات السيولة:} تخصيص الأصول المعتمدة للأغراض المعلنة لمجموعة الخدمات أو خطوط الخدمة تحت الحد الأقصى ومراجعة غروب الشمس.
\item \textbf{التحكم الخارجي في العبور:} التمويل فقط عندما يتم الوفاء بمتطلباتها المنشأة المنشورة وعلى أعلى مستويات الأولوية؛ وإلا تخصيص الصفر.
\item \textbf{الميزانية المقترحة للحسابات:} تخصيص أي بقية مُعتمدة لمجموعات المحافظة على الرسوم المحددة ونشر \texttt{F\_epoch}، والتي قد تكون صفر.
\end{enumerate}
يمكن أن تنظر القرارات الميزانية في الوفاء المؤكد، والمخاطر المغطاة، والاحتياطيات المؤهلة، والحد من الاستخدام، والطرق التي يتم تنفيذها، والتركيز، والحوادث، وأداء الضامن. تتبع كل مقياس قواعد الأدلة والجماعة في الملحق C.

تدفق المساهمين المحتمل سيكون:

\begin{enumerate}[leftmargin=*]
\item يقوم المساهمون بنقل الأصول المؤهلة بموجب شروط المنحة أو البرنامج المعتمدة بشكل منفصل.
\item يحصل المشاركون المؤهلون على رموز إدارة CLC المقترحة، وإذا سمح لهم بذلك، ويقفلونها للحصول على stCLC؛
\item تُدخل شلالات المياه إيصالات ريك أو رسوم خدمة قابلة للتأهيل.
\item اعتمدت صناديق آبشار الأولويات بالترتيب؛ و
\item يتم إصدار sCLC فقط من خلال سياسة ما بعد الاضطرابات المائية المفعولة أو يفتح نافذة محدودة للائتمان من الرسوم.
\end{enumerate}
لا توجد خطوة تخلق حقوق ملكية أو سداد أو سحب أو تغطية أو مكافأة أو حوكمة تتجاوز تلك المنصوص عليها صراحة في الشروط المعمول بها.


\section*{8. النطاق التقني والنمو}
\addcontentsline{toc}{section}{8. النطاق التقني والنمو}

يصف هذا الفصل العمل الاختياري أو المقترح. هذه ليست قائمة بالخصائص التي يضمن وجودها في CLC App الحالي أو Protocol v1.1.0.

تشمل مجالات العمل الممكنة:

\begin{itemize}[leftmargin=*]
\item توجيه التنفيذ عبر مجموعات متوافقة، مع تسجيل، اقتباس، الحد، الرسوم، واكتشاف المخزون؛
\item أجهزة تعديل الاحتفاظ بالتوقيت أو HTLC للتنفيذ عبر النطاقات عندما لا يكون التسوية الذرية متاحة.
\item واجهات وأدوات سياسية للمجمعات الصغيرة أو الشخصية
\item سجلات قابلة للتدقيق للحصول على قسائم، ومجموعات، وأساليب سعر الصرف، والحدات، والمراقبين، والرسوم؛
\item وصلات مزود الدفع الخاصة بالتنفيذ، وتدفقات التسجيلات، ومراقبة الأهلية؛
\item وصلات الأصول المعدلة إلى مراكز السيولة الخارجية لإعادة التوازن وسيولة المدفوعات و
\item تحويل الخزانة المغطاة بالسياسة للتغطية المعتمدة أو تكاليف التشغيل أو تفويضات السيولة.
\end{itemize}
أسعار السوق الخارجية لا تحدد ما يدين به الجهة المصدرة بموجب شروط قسائم. يمكن للمجموعة استخدام مرجع خارجي محمي للأصول المثبتة، ولكن طريقة سعر الصرف المنشورة، والحدود، والرسوم، والكتابة ستحكم أسعارها.

\subsection*{8.1 القواعد المقترحة لخدمة الطرق و SDK}
\addcontentsline{toc}{subsection}{8.1 القواعد المقترحة لخدمة الطرق و SDK}

\textbf{اكتشاف.} ستسأل خدمة الطريق المقترحة عن السجلات المحددة لإدخال الأصول وأساليب سعر الصرف والحدود والرسوم والمخزون والحوادث ومعلومات المراقب. سجلات مخزنة ستشمل حدود الطازجة وتعريفات المصدر.

\textbf{ملفات الشبكة} يمكن للعميل دعم أكثر من أحد أصول السجل أو ملفات السياسة. وسوف يخبر المشارك بالملف، والجهات المقابلة، والمتكيفات، والقيود، ومشغل الخدمات المسؤولة تستخدم الاقتباسات. يجب أن يستوفي مسار متعدد الملفات جميع شروط القفز المطبقة.

\textbf{سياسة المسار} يمكن للمشغل المسؤول استبعاد الاعتمادات غير الآمنة أو الأطراف المقابلة وتطبيق حدود على مستوى المسار ومتطلبات الطازجة ومعايير الصحة. هذه الإشارات ستدعم القرار لن تضمن الوفاء أو الحماية من الخسارة.

\textbf{الرسوم والحدود} سيتم إعطاء اقتباسات عن رسوم المجموعة، وأي رسوم بروتوكول إضافية حالية، وأي رسوم توجيه أو خدمة مقترحة بشكل منفصل. التنفيذ يرفض الاقتباسات التي انتهت صلاحيتها أو الحدود المخالفة.

\textbf{الذرية والتعافي} الإعدام المتعدد القفز سيكون ذري حيثما أمكن عندما استخدمت الـ HTLCs أو الـ escrow، ستفصح الخدمة عن التوقيتات، وسلك المسارات، ومراقبين مسؤولين، وإجراءات الحوادث، والمخاطر المتبقية.

\textbf{المقترح لتحقيق التوازن} قد تجمع خدمة الاختيار عواقب إعادة التوازن والبحث عن دورات أو سلاسل متوافقة. سيكون:

\begin{enumerate}[leftmargin=*]
\item تنشر إيصال يمكن قراءته الآلية يحدد الدورات المفعولة والأصول والكميات وخطوط زمنية التقييم والرسوم؛
\item تنفيذ القيود المعتمدة لفترة واحدة وسياسات الطرف الآخر؛
\item رفض النشاط الذي ينتهك أي ترخيص أو حدود أو مخزون متاح للمجموعة المشاركة. و
\item الحفاظ على المدخلات والإيصالات المحددة للمراجعة ومعالجة النزاعات.
\end{enumerate}
\textbf{متطلبات SDK.} ستوفر SDK للطرق التي يتم تنفيذها خريطة محددة من الاقتباسات إلى الاستلام ، والتحقق من عدم التغيرات لكل جولة ، وكودات الفشل المفهومة ، والسجلات الصديقة للمراجعة. يقدم Protocol v1.1.0 \texttt{SwapRouter} الحالي اقتباسات فقط. لا تقوم بتنفيذ هذه الطرق المقترحة.

\subsubsection*{8.1.1 مواصفات التوافق الحد الأدنى للاتحاد}

إن النظام الإيكولوجي للمجموعة الذي يبحث عن توجيه متعدد الملفات سوف ينشر المعلومات القابلة للقراءة الآلية ل:

\begin{enumerate}[leftmargin=*]
\item \textbf{جذور السجل:} المعرفات للأصول، ومجموعات، وأساليب سعر الصرف، والحدود، وسياسات الرسوم، أو جذور واحدة التي تحل لهم بشكل محدد.
\item \textbf{الإيصالات:} الملف الشخصي، الأصول الداخلية والخارجية، والمبالغ، مصدر الاقتباسات والخاتم الزمنية، الصورة الفورية الحد، الرسوم، نتيجة المخزون، ونتيجة التنفيذ لكل قفزة.
\item \textbf{الإشارات التشغيلية:} معلومات محدودة بالجدد حول المخزون، والحد من الاستخدام، والحوادث، وأي إنجاز أو حماية تمول بشكل منفصل.
\item \textbf{القيود السياسية:} الأطراف المقابلة المسموح بها أو المرفوض بها، فئات الأصول، المعدلات، وأي متطلبات الاحتفاظ بها.
\item \textbf{رموز الفشل:} تفسيرات تحديدية لرفض، انتهاء الصلاحية، الحد، المخزون، السياسة، الاعتماد، أو فشل الحوادث.
\end{enumerate}
يمكن للملف الشخصي إضافة تغطية أو الامتثال أو التحكيم أو خدمات أخرى دون جعلها متطلبات لموافقة CPP الأساسية. ستحدد كل خدمة اختياري طرفها المسؤول وسلطته ونطاقها وشروطها.

\subsection*{8.2 الترخيص والتحقق والخروج}
\addcontentsline{toc}{subsection}{8.2 الترخيص والتحقق والخروج}

عقود Protocol v1.1.0 متوافقة EVM-. يتم نشر العقود في سجل \texttt{src} من مخزن البروتوكول تحت AGPL-3.0 باستثناء المكونات الثالثة غير المعدلة المحددة التي تحتفظ بشروطها الخاصة. يُدعم المصدر والمؤشرات المطبوعة والإرشادات المتعلقة بالتنفيذ مراجعة مستقلة، لكنها لا تثبت بمفردها مراجعة أو تنفيذ آمن أو الامتثال القانوني.

كل عملية تنفيذ سوف تكشف عن نسخة رمزها بشكل منفصل ، ومصدر البناء ، والعناوين ، وسلطات التحكم والترقية ، وحالة المراجعة ، ومرايا السجل ، وأي حماية وقتاً أو توقف.

المقترح \textbf{كيت الشوكة} يمكن أن تشمل:

\begin{enumerate}[leftmargin=*]
\item نصوص النشر المحددة؛
\item صور سريعة للسجل وأدوات التصدير
\item عملية مُوثقة لإعادة توجيه خدمات الطرق وال SDKات والواجهات إلى جذع سجل جديد.
\item قائمة تفتيش مشرف الصندوق لخروج سجل مشترك بأمان؛ و
\item قائمة تفتيش الهجرة للحصول على قسائم لا تنفيذها، بما في ذلك إشعارات الجهة المصدرة وأوقات التقديم والإنجاز، والاستمرار في الوصول إلى السجلات، والإجراءات التعديلية.
\end{enumerate}
يمكن لفرشات متوافقة تحسين الصمود عندما تحتاج المجتمعات والتعاونيات والوكالات العامة والاتحادات والشركات المتعددة أو مشغلي الخدمات إلى إدارة مختلفة. الاستمرارية الفعلية لا تزال تعتمد على ملكية العقد والمفاتيح والاعتمادات والواجهات والبنية التحتية والالتزامات القانونية وخدمات الطرف الثالث.


\section*{9. الاقتصاد المقترح لبرامج السيولة}
\addcontentsline{toc}{section}{9. الاقتصاد المقترح لبرامج السيولة}

يصف هذا الفصل نموذجًا مستقبليًا تم تبنيه بشكل منفصل. إنها ليست ميزة تطبيق حالية أو عرض أو عودة وعد بها أو حق تم إنشاؤه عن طريق الإيداع في \texttt{SwapPool}.

\subsection*{9.1 مصادر الدخل المقترحة}
\addcontentsline{toc}{subsection}{9.1 مصادر الدخل المقترحة}

ميزانية الشبكة المستقبلية يمكن أن تحصل على:

\begin{enumerate}[leftmargin=*]
\item الإفصاح عن \textbf{الشبكة} يتم أخذها كمشاركة في الرسوم المستحقة من المجموعات المشاركة.
\item رسوم التوجيه أو خدمة منفصلة عن الخدمات المشتركة المنفذة؛ و
\item الإيرادات الأخرى المعتمدة صراحة.
\end{enumerate}
الرسوم الإجمالية للمجموعة التي تحتفظ بها المجموعات ليست عائدات الشبكة. تستخدم Protocol v1.1.0 الحالية نموذجًا مختلفًا: يتم إضافة رسوم البروتوكول الاختياري إلى رسوم المجموعة وتُرسل مباشرة إلى المستلم المثبت.

\subsection*{9.2 الرياضيات التوضيحية}
\addcontentsline{toc}{subsection}{9.2 الرياضيات التوضيحية}

إذا كانت الجمعية المشاركة تتقاضى رسوم الجمعية 2.00\% وتتلقى جمعية الشبكة المعتمدة 20\% من هذه رسوم الجمعية، فإن معدل جمعية الشبكة الفعلي المقترح على القيمة المتوجهة هو:

\texttt{rake\_rate = 2.00\% $\times$ 20\% = 0.40\% = 40 bps}

إذا كان 25٪ من الأصول التي تم استلامها والرسوم الخدمية مؤهلة و قابلة للتحويل للاستخدام المعلن بالعملات النقدية بعد التكاليف، فإن الحصة القابلة للاستخدام بالعملات النقدية لـ 40 نقطة قياسية هي حوالي 10 نقطة قياسية.

إيرادات الشبكة المقترحة هي:

\texttt{network\_rake\_received + routing\_or\_service\_fees\_received}

لا تضيف رسوم المجموعة الإجمالية إلى ريك الشبكة: الريك هو تحويل من تلك الرسوم وإلا سيتم احتسابها مرتين.

\subsection*{9.3 حقوق برنامج السيولة}
\addcontentsline{toc}{subsection}{9.3 حقوق برنامج السيولة}

يمكن لبرنامج اعتمد بشكل منفصل تمويل مخزون المجموعة، وخدمات التوجيه، المراقبة، أو ولايات أخرى. يجب أن تكشف شروطها:

\begin{itemize}[leftmargin=*]
\item ما إذا كان التحويل هو هدية أو منح أو قرض أو مساهمة قابلة للإسترداد أو شراء؛
\item الوصاية والتحكم؛
\item قواعد السحب والسداد والخسائر والأولوية؛
\item تأهل الرسوم والمكافأة؛
\item حقوق الحوكمة
\item أساليب تقييم وتقديم التقارير؛ و
\item الإيقاف والإيقاف والإصلاحات
\end{itemize}
لا تخلق \texttt{SwapPool} الحالية أي رموز حصة الجمع أو حقوق المساهمين التلقائيين. يجب أن تستند أي مقياسات "ما بعد" إلى الإيرادات والخسائر المحققة، ولا يجب أن يتم تقديمها كإيرادات مُوعدة، ويمكن أن تكون صفرية أو سلبية.


\section*{10. إطار شامل للمخاطر}
\addcontentsline{toc}{section}{10. إطار شامل للمخاطر}

هذا الإطار المقترح يفصل بين عشرة فئات مخاطر. وتحدد لكل فئة مؤشرات محتملة ومراقبة، واختبارات الإجهاد، وشهية المخاطر الإرشادية. هذه توصيات التصميم، وليس دعوى بأن كل التحكم يتم تنفيذها، أو فعالة، أو كافية. الحدود والاحتياطيات والضمانات والمراقبة والغطاء والحوكمة لا يمكن القضاء على الخسائر.

\subsection*{10.1 مخاطر البروتوكول والعقود الذكية}
\addcontentsline{toc}{subsection}{10.1 مخاطر البروتوكول والعقود الذكية}

\begin{itemize}[leftmargin=*]
\item \textbf{تهديدات:} أخطاء العقود، أخطاء الترقية، فشل الإعتمادات، والحدات أو الرسوم غير المثبتة.
\item \textbf{المؤشرات:} نتائج المراجعة، وتحركات المخزون غير المفسرة، والفشل غير المتغير، والعكس غير العادي.
\item \textbf{التحكم المحتمل:} المراجعات المستقلة؛ الحد الأدنى من الأدوار المميزة؛ المراقبين المكلفين بالوكالة والاعتماد على الإفصاح عنهم؛ التحديثات المحددة في الوقت المناسب؛ المراقبة على السلسلة توقف الحوادث مع السلطة والمعايير المنشورة؛ وطريق هجرة اختبر
\item \textbf{اختبارات الإجهاد:} تعتمد الأسعار غير المتاحة أو الحد، نقص المخزون، تعليق العقود، وتفجير حركة المرور.
\item \textbf{شهية المخاطر:} منخفضة قبل تحديد حجم الالتزامات المتبقية أو حجم التبادلات.
\end{itemize}
\subsection*{10.2 المخاطر الاقتصادية والسوقية}
\addcontentsline{toc}{subsection}{10.2 المخاطر الاقتصادية والسوقية}

\begin{itemize}[leftmargin=*]
\item \textbf{تهديدات:} مخزون رقيق، التدفقات الواحدة الجانبية، السحب السريع، ومراجع الأسعار الملائمة أو القديمة.
\item \textbf{المؤشرات:} الاستفادة العالية من ميزان العملات الرقمية، وتوسيع اختلافات الاقتباسات، والرفض المتكرر للحد، والمخزون المركز.
\item \textbf{التحكم المحتمل:} السقف الحالي لموازنة الرمزيات في المجموعة القيود المقترحة للعمل أو الحسابات؛ الاحتياطيات المعتمدة بشكل منفصل إشارات الأسعار المحمية استثناءات الطريق و رسوم أو حدود الحوادث المحدودة في الوقت.
\item \textbf{اختبارات الإجهاد:} التحركات الكبيرة في الأسعار المرجعية، وارتفاعات العرض، وطلبات السحب، وقطع مصادر البيانات.
\item \textbf{شهية المخاطر:} المعتدلة فقط ضمن المعلمات المنشورة والقدرة على تحمل الخسائر الممولة.
\end{itemize}
\subsection*{10.3 مخاطر الجهة المصدرة والبطاقة}
\addcontentsline{toc}{subsection}{10.3 مخاطر الجهة المصدرة والبطاقة}

\begin{itemize}[leftmargin=*]
\item \textbf{تهديدات:} الإصدار الذي يتجاوز قدرة الوفاء، وعدم الوفاء بالجهة المصدرة، وشروط مضللة، ونوافذ التوافر أو العرض غير المحددة بشكل جيد.
\item \textbf{المؤشرات:} انخفاض معدلات الوفاء المؤكد، وتقديم البطاقات المتعلقة بالعمر، والشكاوى التي لم يتم حلها، والتعرض المركز لمصدر واحد.
\item \textbf{التحكم المحتمل:} التدقيق الواجب مع الجهة المصدرة قسيمة واضحة وشروط العرض حدود الإصدار أو القبول؛ السندات أو الضمانات الممولة بشكل مستقل. الإبلاغ عن المجموعة؛ ومراجعة السجل المسؤولة.
\item \textbf{اختبارات الإجهاد:} الإفلاس للمصدر، والصدمات الإقليمية للإنتاج، والطالبات المزيفة، والتأخيرات الطويلة في الوفاء.
\item \textbf{شهية المخاطر:} انخفاض مع زيادة تركيز الجهة المصدرة أو العرض.
\end{itemize}
\subsection*{10.4 مخاطر تقديم القسيمة للاسترداد والوفاء}
\addcontentsline{toc}{subsection}{10.4 مخاطر تقديم القسيمة للاسترداد والوفاء}

\begin{itemize}[leftmargin=*]
\item \textbf{تهديدات:} الوصول غير صالح أو مزدوج، وعدم كفاية قدرة الجهة المصدرة، وتخفيضات المخزون، وفشل في الخدمات اللوجستية، ونقص سجلات الإبراء.
\item \textbf{المؤشرات:} تأخير طلب الوفاء، وتقديمات فشلت أو مثيرة للجدل، وتخفيضات المخزون، وتخفيضات التذاكر، والوحدات المكتملة التي تفتقر إلى الإبراء.
\item \textbf{التحكم المحتمل:} الإجراءات المنشورة لتقديم المعلومات والتنفيذ؛ الإفصاحات عن القدرة؛ أماكن الوفاء المتعددة حيث يكون ذلك مشروعًا. معايير الأدلة؛ مسارات الشكاوى والعلاج؛ والسجلات التي تفصل بين التقديم والوفاء والإبراء
\item \textbf{اختبارات الإجهاد:} حجم العرض مرتين إلى أربعة أضعاف، وانقطاع المرافق، فشل الموردين، ومحاولات الاستخدام المكرر.
\item \textbf{شهية المخاطر:} منخفضة حيث تتضمن السلع الأساسية أو المشاركين الضعفاء أو نوافذ التوفير الطويلة.
\end{itemize}
\subsection*{10.5 مخاطر الحوكمة}
\addcontentsline{toc}{subsection}{10.5 مخاطر الحوكمة}

\begin{itemize}[leftmargin=*]
\item \textbf{تهديدات:} التقاط المدير أو المراقب، تغييرات المعلمات السريعة، تضارب المصالح، القوى التقنية الخفية، والنداءات الضعيفة.
\item \textbf{المؤشرات:} السلطة المركزة، والإجراءات الطارئة المتكررة، والتغييرات غير المفسرة في السياسات، والتجاوزات المتكررة.
\item \textbf{التحكم المحتمل:} الإفصاحات عن الدور والقوة؛ حدود الموافقة النسبية؛ الأوقات الزمنية الصراعات وقواعد الرفض؛ سجلات التغيير العامة؛ الاستئناف؛ غروبات شمسية ذات طاقة الطوارئ الآلية. و القدرة على الانكشاف
\item \textbf{اختبارات الإجهاد:} المقترحات المتناقضة، وفقدان الموقعين، محاولات الرشوة، والقبض من خلال التراكم أو السيطرة المفوضية.
\item \textbf{شهية المخاطر:} منخفضة في الإجراءات التي تؤثر على أساليب القيمة أو السحب أو جذور السجل أو التغطية أو سلطات الطوارئ.
\end{itemize}
لتنفيذ مستقبل رمز الحوكمة المقترح CLC، سيتضمن اختبار الاستيلاء على التراكم الذي يليه محاولات إعادة توجيه الميزانيات أو إضعاف معايير السجل أو الموافقة على ولايات الأطراف ذات الصلة أو استنزاف التغطية الممولة. وتشمل الضمانات المحتملة إيقافات الحوكمة، وفرق أعلى للأفعال الحرجة، وتأخير التنفيذ، ومراقبة التفويضات والتركيز الشفافة، وعملية الحوادث، وإجراء موثوق بالفرق والخروج. لا يتم تمثيل أي منها على أنه تم نشره ببساطة عن طريق الظهور هنا.

\subsection*{10.6 المخاطر القانونية والامتثال}
\addcontentsline{toc}{subsection}{10.6 المخاطر القانونية والامتثال}

\begin{itemize}[leftmargin=*]
\item \textbf{تهديدات:} قسيمة أو خدمة أو ترقية أو أصول حوكمة تتلقى معاملة تنظيمية غير متوقعة؛ فشل في حماية المستهلكين غسيل الأموال أو التعرض للعقوبات؛ والقيود عبر الحدود.
\item \textbf{المؤشرات:} علامات الولاية القضائية، استفسارات الجهات التنظيمية، الشكاوى، تطابقات الأطراف المقيدة، والانفصال بين السلوك الإعلاني والفعلي.
\item \textbf{التحكم المحتمل:} المراجعة حسب الفئة والولاية القضائية؛ الإفصاحات الدقيقة؛ الواجهات الجغرافية عمليات التحقق المتناسبة من المؤهلات أو التحقق من الشهادات؛ مراقبة الترويج؛ سجلات السلطة والقبول؛ والأطراف المسؤولة واضحة
\item \textbf{اختبارات الإجهاد:} القيود القضائية، وإنهاء المزود، وإعادة التصنيف الإلزامي، وأمر بإيقاف ميزة أو فئة أصول.
\item \textbf{شهية المخاطر:} منخفضة تقييد أو وقف النشاط غير المدعوم.
\end{itemize}
\subsubsection*{10.6.1 الموقع القانوني ومعاملة الرموز المقترحة}

\begin{enumerate}[leftmargin=*]
\item \textbf{البنية التحتية التي يمكن التحقق منها:} تعتبر عقود Protocol v1.1.0 متوافقة مع EVM- ويتم نشر مصدرها بموجب الترخيصات والاستثناءات التي تم تحديدها في مخزن البروتوكول. إن النشر يسمح بالمراجعة لكنه لا يثبت نفسه مراجعة، أو نشر آمن، أو الامتثال القانوني. يجب على كل عملية نشر الكشف عن نسخة رمزها ومصدر البناء والعناوين وسلطات المراقب وحالة المراجعة.
\item \textbf{وضعية الرمز المقترحة:} في هذا التصميم، فإن رمزا حوكمة CLC المقترحة من شأنها تنسيق الحوكمة والوصول المحدد للسياسات. فإنه لن يخلق أرباحا، أو تقاسم الأرباح، أو حقوق بقية، أو ضمان قيمة أو سيولة.
\item \textbf{الأصول المقترحة ذات الصلة} سياسة تنفيذية منفصلة قد تسمح بحجب رمز الحوكمة CLC المقترح لـ stCLC وقد تسمح بتصوير sCLC على مستوى العصر. يجب أن تحدد الشروط المعتمدة حقوقها الدقيقة، وحدودها، انتهاء صلاحيتها، قابلية التحويل، والمعاملة.
\item \textbf{الاتصالات:} لا ينبغي أن تعد المواد لأي إشارة إدارة CLC المقترحة، أو stCLC، أو sCLC، عن الربح أو التقدير أو الدخل السلبي، أو الوصول المضمون.
\item \textbf{مراقبة الولاية القضائية:} قد يتطلب تنفيذ واجهة الجغرافيا، أو شهادات للفئات المقيدة، أو حدود الترويج، أو مراجعة خاصة للأصول، أو التحكمات التي تعيق ميزة مقترحة.
\item \textbf{إشعار المشاركين:} وسوف تشرح الشروط المعتمدة متى يمكن تقليص أو تعطيل الوصول لأسباب قانونية أو تشغيلية أو مخاطر، وما إذا كانت أي تعويضات أو تعويضات تنطبق.
\end{enumerate}
\subsection*{10.7 مخاطر التوجيه وترابط المجال}
\addcontentsline{toc}{subsection}{10.7 مخاطر التوجيه وترابط المجال}

\begin{itemize}[leftmargin=*]
\item \textbf{تهديدات:} التنفيذ الجزئي، والانفجارات المتعثرة، واستغلالات الجسر أو الاحتفاظ، والاقتباسات القديمة، والتضخم في المسار، والقيام بالتحركات المسبقة للتغيرات القيمة المعلنة.
\item \textbf{المؤشرات:} معدلات انتهاء الصلاحية، وتخلفات الاحتفاظ بها، وفروق بين الاقتباسات إلى التنفيذ، وارتفاعات غير ضرورية متكررة، وحوادث الجسر.
\item \textbf{التحكم المحتمل:} التنفيذ الذري إذا كان متاحًا. الاحتياطيات HTLC أو الاحتياطيات الاحتياطية سياسات المسار والجهات المعارضة؛ الحواجز على مستوى المسار خريطة تحديدية من الاقتباسات إلى الاستلامات وشركات الخدمات المسؤولة.
\item \textbf{اختبارات الإجهاد:} توقف جسر، إعادة تنظيم سلسلة، انقطاع الاعتماد، و قفزة واحدة فشلت في مسار متعددة القفزات المقترحة.
\item \textbf{شهية المخاطر:} منخفضة إلى معتدلة فقط بالنسبة للاعتمادات المحددة والمراقبة.
\end{itemize}
\subsection*{10.8 مخاطر الحماية وإدارة المفاتيح}
\addcontentsline{toc}{subsection}{10.8 مخاطر الحماية وإدارة المفاتيح}

\begin{itemize}[leftmargin=*]
\item \textbf{تهديدات:} المفاتيح المفقودة أو المعرضة للخطر، التواطؤ مع الموقّع، والتحميل غير الواضح أو سلطة المدير.
\item \textbf{المؤشرات:} توقيعات غير عادية، تغييرات في جهاز التحكم، فشل في الدوران، والسحب غير العادي.
\item \textbf{التحكم المحتمل:} الترخيص المتعدد الأجزاء أو الحد الأدنى؛ المفاتيح المدعومة بالأجهزة؛ الفصل بين الأدوار تدوير مؤشر؛ مخزونات المراقبين العامين حدود السحب المراقبة؛ وإجراءات الاسترداد التي تم اختبارها.
\item \textbf{شهية المخاطر:} منخفضة
\end{itemize}
\subsection*{10.9 السمعة والخطر الاجتماعي}
\addcontentsline{toc}{subsection}{10.9 السمعة والخطر الاجتماعي}

\begin{itemize}[leftmargin=*]
\item \textbf{تهديدات:} الادعاءات المضللة، الحوافز الضارة، الشكاوى التي لا يمكن الوصول إليها، الخبرات السيئة في الوفاء، فشل الخصوصية، والحماية التي تفضي إلى الداخلين.
\item \textbf{المؤشرات:} الشكاوى حسب المجموعة، والنزاعات التي لم يتم حلها، والاتجاهات في أداء الجهة المصدرة، وتركيز الفوائد أو الخسائر، والردود المجتمعية.
\item \textbf{التحكم المحتمل:} الإفصاحات باللغة البسيطة؛ مسارات الشكاوى والتصحيح الأدلة المتاحة؛ الإبلاغ الشفاف؛ مراجعة الحوادث والعقوبات المتناسبة على الإدعاءات الخاطئة.
\item \textbf{شهية المخاطر:} منخفضة، مع رعاية خاصة للمجتمعات المتضررة والمشاركين الضعفاء.
\end{itemize}
\subsection*{10.10 مخاطر التركيز والتجزئة}
\addcontentsline{toc}{subsection}{10.10 مخاطر التركيز والتجزئة}

\begin{itemize}[leftmargin=*]
\item \textbf{تهديدات:} الاعتماد على عدد قليل من الجهات المصدرة، ومجموعات، ومراقبين، ومقدمي الخدمات، وشبكات، أو شوك غير متوافقة.
\item \textbf{المؤشرات:} تدابير التركيز من قبل الجهة المصدرة أو المجموعة أو المخزون أو المراقب أو مزود الخدمات؛ أخطاء الطريق بين المجموعات؛ والاعتمادات الوحيدة الحرجة.
\item \textbf{التحكم المحتمل:} حدود التركيز المنشورة؛ العديد من المشغلين المسؤولين؛ المعايير المتوافقة السجلات المستقلة إجراءات الخروج المختبرة و التفاعلية الآمنة.
\item \textbf{اختبارات الإجهاد:} فقدان أكبر مصدر أو مجموعة أو مشغل أو مزود أو جذور السجل.
\item \textbf{شهية المخاطر:} خاصة بالتنفيذ والكشف عنها، مع قيود أكثر صرامة للخدمات الأساسية أو الاعتمادات التي لا يمكن استبدالها.
\end{itemize}

\section*{11. آليات الحوكمة}
\addcontentsline{toc}{section}{11. آليات الحوكمة}

يقترح هذا الفصل نموذج حكم. هذا لا يمثل أن CLC App الحالية تستخدم التصويت على رموز الحوكمة، وقاعات التوقيت، التأمين المشترك، عملية المطالبات، أو كل التحكم الموصوف أدناه. كل عملية تنفيذ تحتاج إلى تحديد صانعي القرار الفعليين والسلطات والعقود والعمليات والسياسات.

\begin{itemize}[leftmargin=*]
\item \textbf{القيم الدستورية:} رعاية الناس، رعاية البيئة، العدالة، المتبادلة، عدم الهيمنة، والمرونة.
\item \textbf{أنواع المقترحات:} تغييرات في الرسوم والحد والمعايير؛ ولايات السيولة؛ إدراجات المجموعة وإزالتها قرارات التغطية الاختيارية والحواجز المحددة.
\item \textbf{العملية المسؤولة:} الإدخال $\to$ التقييم $\to$ مراجعة المخاطر $\to$ الموافقة $\to$ التوقيت عند الاقتضاء $\to$ التنفيذ. يمكن أن تأتي الموافقة من المديرين أو التعاونيات أو الوكالات العامة أو الاتحادات أو الشركات المتعددة أو التصويت على السلسلة أو من هيئة أخرى مطلعة ومسؤولة.
\item \textbf{حدود الموافقة:} المعلمات حسب فئة الإجراءات، مع أعلى حدود للتغييرات في مؤشر القيمة، والسلطات الطارئة، وغيرها من الإجراءات الحرجة.
\item \textbf{الوفد:} التفويض الاختياري مع تفويضات عامة، الإفصاح عن الصراعات، والاستدعاء.
\item \textbf{قطع الدوائر:} وقفات الطوارئ مع المعايير المعلنة، والعملاء المعتمدين، وشروط الاستئناف، والتحقيقات المطلوبة بعد الوفاة.
\item \textbf{الشفافية:} تغييرات وتدفقات نشرت، مع أدلة منفصلة على تسوية المبادلات المبادلة، وفقا للمصدر، والاحتياطيات، واستخدام الحد، التوجيه، والضامنين.
\end{itemize}
\textbf{إدارة السجلات} يمكن لنشر CPP- المتوافق أن يحافظ على سجلات اكتشاف البطاقات والرموز والسباحات. يمكن للضوابط المعتمدة إضافة أو تحديث أو تعليق أو إزالة إدخالات السجل من خلال عملية الحوكمة المعلنة للتنفيذ. إزالة السجل يؤثر على اكتشاف وتوجيه عبر هذا السجل. فإنه لا يمحو بطبيعة الحال رمز، أو يغير رصيد حامل، أو يفي بالتزامات الجهة المصدرة، أو يمنع عقدًا وظيفيًا آخر.

وينبغي أن تكون قواعد السجل المنشورة مشروطة على الوضع، ويمكن أن تحدد عدم الوفاء المتكرر أو الاحتيال أو التعبير الخاطئ أو سلوك عقد غير آمن أو انتهاك مستمر للمبادئ المنشورة كأسباب للايقاف أو الإزالة. حيثما كان ذلك ممكناً، يجب أن توفر العملية إشعاراً وفرصة للتعويض وطرق الاستئناف. يجب أن يتطلب إزالة الطوارئ تقرير حادثة عامة ومراجعة تلقائية أو غروب الشمس.

\textbf{القوائم المحظورة} في هذا النموذج، لا يسمح السجل:

\begin{enumerate}[leftmargin=*]
\item الأدوات التي تمول أو تحفز بشكل مباشر التدمير البيئي خارج الحدود المتفق عليها، والعنف أو استخدام الأسلحة، والاستخراج القسري، أو الإساءة النظامية؛ أو
\item فئة البطاقات التي تفتقر إلى شروط واضحة لتقديم وتوفير، والمساءلة، وطرق العلاج.
\end{enumerate}
ستكون القائمة المحظورة نسخة، قابلة للتدقيق العام، ويمكن تغييرها فقط من خلال العد الأدنى المعتمد للعمل الحرج والقيود الزمني، كما هو موضح على شكل Q3 + T3 في الملحق D.

\subsection*{11.1 سعر الصرف وإدارة الحدود}
\addcontentsline{toc}{subsection}{11.1 سعر الصرف وإدارة الحدود}

\textbf{تغييرات في الوقت المحدد} إن التنفيذ الذي يتبع هذا الشكل سيغير أساليب أسعار الصرف ويتم تنفيذ معايير الحد بشكل منفصل فقط بعد إيقاف الوقت العام. سيقوم مسار الطوارئ باستخدام عملية تفويض مطولة بشكل منفصل وتشمل غروب الشمس أو مراجعة تلقائية.

\textbf{حدود الموافقة} يقترح النموذج أعلى حدود الموافقة لتغييرات قاعدة مؤشر القيمة والتغييرات العالمية على مستوى الحد، والعوائق المتوسطة لتغييرات الطرف الثالث المحددة للمجموعة، والعوائق القياسية لتغييرات الرسوم الروتينية.

\textbf{المواد المنشورة.} سيتم نشر المنشأة المشاركة ، لكل مجموعة ، متغيرات مؤشر السلسلة ، ومصادر أو المتوسطات ، وتحديث الترتيبات ، وطرق الحد من النوافذ والقواطع ، وأوضاع الفشل أو ثابتات آمنة.

\textbf{معايير توقف الطوارئ} سيتم الإعلان مسبقاً عن ظروف مثل انقطاع الكلام، واستخدام الحد المرتفع جنباً إلى جنب مع فشل الإنجاز، أو فشل غير متغير، جنباً إلى جنب مع عمليات التحقق من الاستئناف ومتطلبات مراجعة ما بعد الحادث.

\subsection*{مثال على تغذية المؤشرات العامة لمجموعة واحدة والبطاقة}
\addcontentsline{toc}{subsection}{مثال على تغذية المؤشرات العامة لمجموعة واحدة والبطاقة}

\begin{itemize}[leftmargin=*]
\item \textbf{الرمز:} على سبيل المثال، \texttt{Maize\_50kg@IssuerY}.
\item \textbf{وحدة مرجعية:} وحدة المؤشر (IUX).
\item \textbf{القيمة المنشورة:} 30.000 IUX.
\item \textbf{المصدر:} المتوسط من المصادر المحددة، مثل استطلاع السوق المحلي، والرسالة الوزارية، والخط الأساسي للتنفيذ.
\item \textbf{تحديث التسلسل:} يوميًا في الساعة 18:00 EAT، مع حجب زمني لمدة 24 ساعة.
\item \textbf{وضع الفشل:} التجمد عند آخر قيمة صالحة، وتطبيق سياسة الحد المعلن، والتوقف بعد توقف 72 ساعة.
\item \textbf{الأسباب:} الملاحظات المنشورة وتسجيل التغييرات من التحديث السابق.
\item \textbf{الموقعون:} العناوين المتعددة الإشارات التي تم الكشف عنها وفرق الموافقة.
\end{itemize}
\subsection*{11.2 دليل التدريب المقترح على صناديق التأمين}
\addcontentsline{toc}{subsection}{11.2 دليل التدريب المقترح على صناديق التأمين}

\textbf{التصميم الاختياري فقط} ينطبق هذا الدليل فقط على التنفيذ الذي اعتمد وتمويلا صريحا لصندوق التأمين ونشر الأحداث المغطاة والمطالبين المؤهلين، والكيان المسؤول، والأصول، والحدود، والإقصاءات، ومتطلبات الأدلة، والعملية، وشروط الحكم. لا توفر CLC App الحالية ولا GEF تغطية فقط لأن هذا التصميم ظهر في الورقة البيضاء.

\textbf{محفزات محتملة} يمكن أن تغطي السياسة المعتمدة عدم الوفاء بالجهة المصدرة المحدد، أو نقص في احتياطي المجموعة، أو خسارة في الجسر أو الاحتياطي. الحوادث التقنية لا تصبح مؤهلة تلقائيًا. السياسة المعمول بها سوف تتحكم.

\textbf{تقييم} ستقوم الهيئة المسؤولة بتوفيق استلامات المعاملات، رصيدات المخزونات، سندات الضامن، سجلات تقديم القسيمة للاسترداد، استجابات الجهة المصدرة، وغيرها من الأدلة المطلوبة، ثم تنشر سجل الحوادث بما يتوافق مع الخصوصية والقانون.

\textbf{شلال خسارة مثالي} عندما توجد كل طبقة وتطبق قانونياً، يمكن أن تستخدم السياسة: (1) سندات الجهة المصدرة المسؤول أو أسهم الضامن $\to$ (2) احتياطيات على مستوى الحجم $\to$ (3) صندوق تأمين الشبكة المقترح $\to$ (4) تخفيض مؤقت لمطالبة التغطية الاختياري، فقط إذا سمحت الشروط الحالية السابقة والقانون المعمول بها صراحة $\to$ (5) استرداد قانوني لإثبات الاحتيال أو سوء الاستخدام.

لا يقلل تعديل التغطية من الالتزام الأساسي للمصدر بالبطاقة أو يغير رصيد السلسلة ما لم تسمح هذه النتيجة صراحةً بالشروط الحالية السابقة والقانون المعمول بها، ويتم الحصول على أي موافقة مطلوبة من حاملها.

\textbf{الحدود والإقصاءات} وتحدد التغطية المنشورة القيود، والتقديمات المؤهلة، والأدلة، ونوافذ المطالبات، والطرق أو الأحداث المستبعدة، والقيود الجغرافية، ومعاملة الاحتياطيات المستنفدة. يمكن أن يكون الدفع صفر بعد الوصول إلى الحدود المعمول بها.

\textbf{جدول إعادة التأهيل الموضح} إذا تم اعتمادها ونشرها:

\begin{enumerate}[leftmargin=*]
\item الادعاءات ستستمد أولاً من الجهة المصدرة المسؤول أو الضامن السندات، ثم من احتياطيات المجموعة المعمول بها، ثم من صندوق التأمين على الشبكة المقترح؛
\item أي تخفيض لمطالبة تغطية اختيارية سيكون محدودًا بما تسمح به شروط تغطية سابقة والقانون المعمول به ، وحتى الحد الأقصى للحوادث المنشورة.
\item يمكن أن تطبق خطة الاسترداد حصة معينة من القيمة المستردة لفترة معينة ، وبعد ذلك فإن أي نقص مغطى متبقي سيصبح خسارة مسجلة مع حسابات عامة بعد الوفاة ؛ و
\item كل قرار يقدم إيصالًا مع الحادث ID، والطلبات والبطاقات المتضررة، والقرار، وخطة التعافي، ونقطة الاستئناف.
\end{enumerate}
\subsection*{11.3 إطار الضمان}
\addcontentsline{toc}{subsection}{11.3 إطار الضمان}

هذا القسم يميز بين مسؤولية الجهة المصدرة وحماية المجموعة الاختيارية وضمانات الطرف الثالث. يمكن للمجمعات التنافس على الحماية والشروط والحماية المقدمة صراحة دون أن يشير إلى أن CLC App، CPP، GEF، أو أي شبكة أوسع تضمن بطاقة بطاقة تلقائيًا.

\subsection*{مسؤولية الجهة المصدرة الأساسي}
\addcontentsline{toc}{subsection}{مسؤولية الجهة المصدرة الأساسي}

\begin{itemize}[leftmargin=*]
\item كل قسيمة هي في المقام الأول مسؤولية مصدرها. يلتزم الجهة المصدرة بتقديم السلع أو الخدمات المعلنة أو ما يعادل النقود القانونية بموجب شروطها المنشورة.
\item يجب على الجهات المصدرة نشر من يمكنهم تقديم البطاقة، وما هو معنى الوفاء بها، أين ومتى تكون متاحة، ما هي الأدلة المطلوبة، وما هي العلاجات المطبقة.
\item إذا فشل أحد الجهة المصدرة في الوفاء، فإن الجهة المصدرة هو الجهة المسؤولة الأساسية. لا تنطبق حماية المجموعة أو الشبكة إلا عندما يتم اعتمادها بشكل منفصل وتمويلها وإفصاحها.
\end{itemize}
\subsection*{حماية البحيرة الاختيارية}
\addcontentsline{toc}{subsection}{حماية البحيرة الاختيارية}

قد يختار مشرف الصندوق إضافة حماية محددة ضيقة إلى البطاقات المقبولة. إنه ليس تلقائيًا ويحتاج إلى تحديد الجهة المسؤولة، والتمويل، والأحداث المؤهلة، والسقفات، والنوافذ، والأدلة، والإقصاءات، والإصلاحات في البيانات الأساسية للمجموعة والشروط المعمول بها.

تشمل أنواع الحماية الموضحة:

\begin{enumerate}[leftmargin=*]
\item \textbf{تغطية الأصول الاحتياطية} بعد عدم الوفاء بالجهة المصدرة المحقق، يدفع الكيان المسؤول في المجموعة مبلغًا محددًا في أصل احتياطي معين، مع مراعاة الحد الأقصى المنشور والاحتياطيات الممولة المتاحة.
\item \textbf{نافذة إعادة التغيير:} بعد حدث مؤهل، يقدم المجموعة مسار تبادل محدود في الوقت إلى الأصول المعتمدة سابقاً أو غيرها، مع مراعاة القيود والتخزين. هذه حماية سيولة تعتمد على المخزون، وليس وعد بأن كل تبادل قابل للتعديل.
\item \textbf{الوفاء البديل:} ينظم الطرف المسؤول مزودًا بديلًا معتمدًا ضمن حد الكمية أو القيمة المنشورة.
\item \textbf{حماية بنطاق الصرف:} بالنسبة إلى فئات البطاقات المختارة، يقدم المجموعة فقط تعديل التغطية أو تعويضات المقايضة المذكورة في شروطها الحالية مسبقاً. هذا لا يقلل من الالتزام الأساسي للمصدر بالبطاقة.
\end{enumerate}
\subsection*{مصادر التمويل الممكنة}
\addcontentsline{toc}{subsection}{مصادر التمويل الممكنة}

\begin{itemize}[leftmargin=*]
\item \textbf{سندات الجهة المصدرة:} الضمانات التي يضعها الجهة المصدرة أو يحتفظ بها في احتياطية مفضحة وتتوفر بعد حدث مُغطى مُحقق.
\item \textbf{احتياطي الصندوق:} الأصول التي تسيطر عليها الكيان المسؤول في المجموعة وتخصص للحماية التي يعلن عنها.
\item \textbf{سندات الضامن من طرف ثالث:} الضمانات التي يضعها ضامن خارجي معين للمصدرين المعلنين أو فئات البطاقات أو الأحداث.
\end{itemize}
سيتبع مشاركة الضامن معايير الأهلية المنشورة، وتحديد حجم السندات، وحدود التركيز، سلطة القرار، وقواعد التنفيذ القانوني.

\subsection*{عملية المطالبات}
\addcontentsline{toc}{subsection}{عملية المطالبات}

ستحدد السياسة المعتمدة العوامل المحفزة للدراسة، مثل الموعد النهائي للقيام بالقيام بالقيام بالقيام بالقيام بعد تقديم الرداءات الصالحة، والإفلات الجهة المصدرة المحقق، والجسر المغطى أو فشل الاحتفاظ، أو حالة الحوادث المعلنة رسميا. سوف تحدد أيضاً:

\begin{itemize}[leftmargin=*]
\item كيف يقوم المشارك بفتح مطالبة وتقديم الأدلة المطلوبة على تقديمها وتوفيرها؛
\item الذي يتحقق من شروط البطاقة، والردود التي يقدمها الجهة المصدرة، والسجلات التقنية؛
\item النوافذ القرارية والإستئنافية و
\item مسار الدفع المسموح به، والأصول، والأسقف، والإيصالات.
\end{itemize}
إن عائدات الاسترداد من الجهات المصدرة أو التحكيم أو التنفيذ القانوني ستقوم بإعادة ملء السندات أو الاحتياطيات المعمول بها وفقاً للسياسة المنشورة قبل استخدامها للوصول إلى تجمع شبكة CLC المقترحة.

\subsection*{الإفصاحات المطلوبة}
\addcontentsline{toc}{subsection}{الإفصاحات المطلوبة}

لكل فئة من المجموعات المغطاة والبطاقات، ينشر الطرف المسؤول:

\begin{itemize}[leftmargin=*]
\item ما إذا كان الضامن غائبًا أو اختياريًا أو مطلوبًا؛
\item حجم السندات أو الاحتياطيات والحد الأقصى لتركيزها
\item أنواع الحماية والأصول والقواطع والنوافذ والاستبعاد؛
\item المواعيد النهائية لتقديم، والوفاء، والمطالبة، والإستئناف؛ و
\item بيان بسيط عن من يضمن ما وما لا يضمن.
\end{itemize}
\textbf{مبدأ الحفاظ} مشرفو الصناديق والهياكل القانونية أو الحوكمة المسؤولة تتحمل المسؤولية عن الحماية التي يعلنون عنها. قد يوفر تنفيذ CPP- متوافق مع المعايير أو السجلات أو السياسات المشتركة الاختيارية، ولكن لا توفر CLC أو GEF بطاقات أو مجموعات تلقائيًا.

\subsection*{11.4 الحواجز المضادة للقبض}
\addcontentsline{toc}{subsection}{11.4 الحواجز المضادة للقبض}

ووفقًا لهذا النموذج، فإن الإجراءات الحرجة التالية تتطلب أعلى مستوى من الموافقة المعتمدة وقتاً طويلاً:

\begin{enumerate}[leftmargin=*]
\item تغيير شلالات الرسوم المقترحة، بما في ذلك تغطيتها وأولويات العمليات الأساسية؛
\item تغيير جذور السجل الكنسي؛
\item تغيير نطاق التغطية أو حدود المطالبات أو سلطة القرار؛
\item توسيع صلاحيات توقف الطوارئ أو
\item إضعاف الالتزامات المتعلقة بالشفافية أو السيادة في المجموعة المذكورة في هذه الورقة.
\end{enumerate}
\subsection*{11.5 الإجراءات المتعلقة بالفوك والخروج}
\addcontentsline{toc}{subsection}{11.5 الإجراءات المتعلقة بالفوك والخروج}

إذا تم الاستيلاء على الحوكمة أو انخفضت القيم بشكل كبير، يمكن للمجتمعات، و مشرفو الصناديق، و المشغلين السعي إلى الخروج عن طريق تشويه طبقة حوكمة الشبكة. استمرارية المجموعات والبطاقات الأساسية تعتمد على العقود المنفذة والمفاتيح والواجهات والبنية التحتية وخدمات الطرف الثالث والالتزامات المعمول بها.

عملية الخروج يمكن أن:

\begin{enumerate}[leftmargin=*]
\item \textbf{نشر اللقطة:} تصدير السجلات المختارة والبطاقات والقيم والحدود وسياسات الرسوم، ثم نشر شريط سريع تم توقيعه.
\item \textbf{إعادة نشر خدمات الحوكمة:} تنشر جذور السجل الجديدة، وخدمات الطرق، وأي وحدات رسوم أو تغطية اعتمدت في إطار هيكل حسابي جديد.
\item \textbf{إعادة التسجيل:} السماح لـ مشرفو الصناديق بالاختيار من خلال تسجيل عناوين المجموعة الخاصة بهم تحت الجذر الجديد دون أن يتطلب من أصحابها نقل قسائم وظيفية أخرى.
\item \textbf{أعد العملاء:} إضافة الجذر الجديد كملف تعريف شبكة قابلة للاختيار في SDKs والواجهات، مع أي تغيير افتراضي يتم من خلال عملية الحوكمة المعلنة.
\item \textbf{إدارة فترة الجسر:} الحفاظ على الطرق المتوافقة حيث تكون آمنة ورفض الطرق التي تنتهك قواعد الملف الجديد.
\end{enumerate}
الهدف من التصميم هو أن ترك السجل الكنسي لا يؤدي إلى تعطيل الصندوق المحلية التي تعمل بشكل آخر. الاستمرارية الفعلية لا تزال تعتمد على الانتشار. الاتحاد هو طبقة الاكتشاف والتنسيق المشتركة.


\section*{12. ورقة زمنية برنامج السيولة المقترحة (مخطط غير ملزم)}
\addcontentsline{toc}{section}{12. ورقة زمنية برنامج السيولة المقترحة (مخطط غير ملزم)}

هذه قائمة تفحص مثالية لبرنامج مستقبلي مُوثق بشكل منفصل ليس عرضًا أو ميزة حالية من التطبيق أو وعداً بالسيولة أو الوصول أو التأمين أو العائد أو القيمة أو الخروج.

\begin{itemize}[leftmargin=*]
\item \textbf{المساهمات المؤهلة:} العملات المستقرة أو رموز الاحتياطي أو بطاقات المجتمع المعتمدة، كما هو موضح في البرنامج.
\item \textbf{الموقع:} تم تعيين صناديق الالتزامات أو مجموعة شبكة CLC المقترحة بموجب ولاية تم اعتمادها بشكل منفصل.
\item \textbf{الحبس والخروج:} الشروط المحددة للبرنامج، مثل فترات 30-90 يوماً وتكشف عن البوابات تحت الضغط.
\item \textbf{ائتمان الرسوم المحددة للسياسة:} آلية اختياريّة في وقتٍ لاحق تحت الحواجز والنوافذ المنشورة، وليس عائدًا مُوعدًا.
\item \textbf{تقاسم المخاطر:} أي تخفيض إلى مطالبة تغطية اختيارية يجب أن يسمح صراحة بموجب الشروط القائمة من قبل والقانون المعمول به. لا تغير نفسها الالتزام بالبطاقة الخاصة بالجهة المصدرة أو الرصيد على السلسلة.
\item \textbf{الإبلاغ} لوحات التحكم الدورية والشهادات التي تغطي صحة المجموعة، والمخزونات، والحدود، والرسوم، وأي احتياطيات التغطية الممولة.
\item \textbf{المعاهدات:} القيود في استخدام الإيرادات، والتحكم في الأوراقل، ومستويات الحد، وقواعد الصراعات، والإجراءات العلاجية.
\item \textbf{الموافقة المحتملة على الرمز المقترح:} يمكن للمساهمين الذين يطلقون على المجموعات المحددة أن يكونوا مؤهلين للحصول على منح رموز حكم CLC المقترحة المخصصة بناءً على الوصول إلى المجموعة المحددة لتبادل المجموعة والإنجاز المؤكد من قبل الجهة المصدرة، مع مراعاة الشروط المعتمدة وقواعد مكافحة الألعاب والسقف السياسية.
\end{itemize}


\section*{13. قاموس لغة بسيطة}
\addcontentsline{toc}{section}{13. قاموس لغة بسيطة}

\begin{itemize}[leftmargin=*]
\item \textbf{Cosmo-Local Credit (CLC):} عائلة المنتجات، بما في ذلك CLC App، والوثائق، والعمل الأوسع لتجميع الالتزامات.
\item \textbf{CLC App:} تطبيق الويب التقدم الذي يديره GEF في \texttt{cosmolocal.credit}.
\item \textbf{بروتوكول تجمع الالتزامات (CPP):} النموذج الذي يمكن إعادة استخدامه من الاحتفاظ به، والتقييم، والحد، والتبادل، والحوكمة المسؤولة.
\item \textbf{Protocol v1.1.0:} الإفراج عن العقود الذكية العامة
\item \textbf{صندوق الالتزامات:} الترتيب الحكم الذي يعترف بالأصول المدعومة و ينشر قواعد الصرف.
\item \textbf{\texttt{SwapPool}:} الصندوق الرمزي الحالي وعقد التبادل المباشر
\item \textbf{الرمز:} رصيد على السلسلة أو الأصول الرقمية آليات الرموز وحدها لا تخلق وعدا قابلا للاسترداد
\item \textbf{البطاقة:} رمز أو سجل تمثله الجهة المصدرة كالتزام قابل للاسترداد بموجب الشروط المنشورة
\item \textbf{العرض:} سجل الكتالوج للسلع أو الخدمات المرتبطة بالبطاقة.
\item \textbf{الجهة المصدرة:} الجهة المسؤولة عن نشر وتحقيق الالتزام بالبطاقة.
\item \textbf{مشرف الصندوق:} الهيكل المسؤول الذي ينشر ويدير قواعد المجموعة.
\item \textbf{سعر صرف المجموعة أو الاقتباس:} معايير المعاملات التي يتم إنتاجها بواسطة مؤشر مدفوع ؛ لا دليل على النقدية أو قيمة الاسترداد.
\item \textbf{الحد الأقصى لتحديد الميزانية:} الحد الأقصى الحالي \texttt{Limiter} للرمز الذي يحتفظ به مجموعة. التطبيق قد يضع عليها علامة \textbf{``حد الائتمان.''}
\item \textbf{تبادل المجموعة:} تبادل أصول مباشر من خلال \texttt{SwapPool}.
\item \textbf{تسوية التبادل:} التحويلات على السلسلة ومحاسبة الرسوم التي تستكمل تبادل المجموعة.
\item \textbf{تقديم القسيمة للاسترداد:} إرجاع أو تقديم وحدات البطاقة إلى الجهة المصدرة.
\item \textbf{الوفاء:} يقوم الجهة المصدرة بتقديم السلعة أو الخدمة أو الفائدة أو الأداء الآخر الموعود به.
\item \textbf{الإبراء:} سجل يمنع عرض الوحدات المكتملة مرة أخرى.
\item \textbf{جهاز توجيه التنفيذ المقترح:} البرمجيات المستقبلية التي من شأنها تفويض وتنفيذ الطرق. الاقتباسات الحالية \texttt{SwapRouter} فقط.
\item \textbf{مجموعة شبكة CLC المقترحة:} ترتيب مستقبلي للتصفية على مستوى الشبكة وتنسيق الميزانية.
\item \textbf{تم اقتراح رمز الحكم CLC:} تصميم إدارة مستقبلي، وليس أسلوب التطبيق الحالي أو Protocol v1.1.0.
\item \textbf{التشويش الشبكي:} حصة مقترحة من رسوم المجموعة التي يتم نقلها إلى ميزانية شبكة مستقبلية.
\item \textbf{ضمان أو تأمين:} الحماية التي تنطبق فقط عندما يتخذ الطرف المحدد صراحة، وتمويلها، ونشرها.
\end{itemize}

\section*{14. المؤشرات الرئيسية والمؤشرات الصحية المقترحة}
\addcontentsline{toc}{section}{14. المؤشرات الرئيسية والمؤشرات الصحية المقترحة}

وينبغي أن ينشر نشر KPI فقط عندما يمكنه تحديد الحدث، المصدر، الجماعة أو الفترة، الوحدة، طريقة التقييم والخيط الزمني، الاستبعادات، سياسة التصحيح، الأدلة خارج السلسلة، والقيود على جودة البيانات.

تشمل الإجراءات المقترحة:

\begin{itemize}[leftmargin=*]
\item تقديمات فدية سارية؛
\item معدل الوفاء القائم على الجماعات؛
\item استكمال الإبراء؛
\item تأخير التقدم إلى الإنجاز؛
\item مدة الاستحواذ من الاستحواذ إلى التقديم
\item الالتزامات المؤهلة المعلنة بموجب منهجية محددة.
\item حجم تبادل المجموعة المكتملة
\item مخزون المجموعة المقياس ؛
\item الاستفادة من رموز الموازنة في المجموعة
\item معدل عبور الاقتراحات ، والذي يتم الاحتفاظ به بشكل منفصل عن معدل تنفيذ الطريق ؛
\item الاحتياطيات المؤهلة مقسمة على التعرض المغطى صراحة.
\item مطالبات الضامن والاسترداد بموجب سياسة محددة؛
\item رسوم الشبكة المقترحة والخدمات التي تلقت فعلياً، دون رسوم المجموعة الإجمالية المزدوجة؛
\item أوقات الكشف عن الحوكمة والقرار والتوقف والإصلاح وإغلاقها.
\item الشكاوى والنزاعات والتصحيحات والإصلاحات و
\item أظهرت بشكل منفصل النتائج الاجتماعية أو البيئية.
\end{itemize}
لا تثبت بيانات المعاملات على السلسلة بذاتها الهوية أو أداء الجهة المصدرة أو الرضا أو السببية أو صحة المجتمع أو التأثير الاجتماعي. هذه الادعاءات تتطلب منهجيتها ودليلها الخاصة.

انظر \href{https://docs.cosmolocal.credit/ar/white-paper/appendix-c-kpi-definitions}{الملحق (ج)} لمواصفات القياس المقترحة.


\section*{15. خريطة الطريق (مؤشرة)}
\addcontentsline{toc}{section}{15. خريطة الطريق (مؤشرة)}

\subsection*{15.1 الأساس الحالي}
\addcontentsline{toc}{subsection}{15.1 الأساس الحالي}

يدير GEF CLC App في \texttt{cosmolocal.credit} على سلسلة Gnosis. يوفر التطبيق إمكانية الوصول إلى الحسابات والمحفظة المدعومة، وكتالوج السوق العام، وواجهات للشعارات، والبطاقات، والعروض، و صناديق الالتزامات، والتحويلات، والتبادلات المباشرة للمجموعة.

توفر Protocol v1.1.0 أساس العقد الحالي: \texttt{GiftableToken}، التنفيذ المباشر \texttt{SwapPool}، السجلات الاختياريّة، وحدات التقييم، أعلى مستويات تكوينات الميزان المجموعة، رسوم المجموعة، رسوم بروتوكول إضافية، و \texttt{SwapRouter} ذات اقتباس فقط. لا يزال التوافر يعتمد على الواجهة والعقود المنشورة والخزون والتنظيم وحالة الشبكة والقدرة على التأهل والمقدمين والولاية القضائية والشروط المنشورة للمصدر أو المجموعة.

\subsection*{15.2 الخطوط الهامة المقترحة}
\addcontentsline{toc}{subsection}{15.2 الخطوط الهامة المقترحة}

هذه الأهداف المشار إليها هي توجيهات التصميم، وليس أرقام الإصدار، أو تواريخ التسليم، أو الالتزامات التي سيتم إطلاق ميزة.

\begin{itemize}[leftmargin=*]
\item \textbf{مؤسسة الحوكمة:} إشارة حوكمة CLC المقترحة تجمع شبكة CLC المقترح أجهزة تعديل الرسوم القرار، وقتا، وأسس الحوكمة.
\item \textbf{التوجيه والتحقق:} الجهاز التوجيه SDK و APIs السجل ؛ أجهزة التحكم الصحية إطار سياسة التأمين المقترح؛ نية إعادة التوازن الاختيارية؛ النموذج الأولي للشبكات
\item \textbf{أدوات المخاطر عبر المجالات:} HTLC أو توجيه الاحتفاظ بها؛ وحدات الضامن التي تحكمها بشكل منفصل؛ إعدادات مسبقة للتدفق، والحساب، والحد المحدد.
\item \textbf{الوصول المنظم:} خدمات الدفع الخاصة بالتنفيذ من طرف ثالث؛ الصناديق الصغيرة الشخصية اكتشاف خدمة الامتثال؛ مراجعات الطرف الثالث لفئات البطاقات.
\end{itemize}
يتم تقديم أي خدمة دفع من قبل أطراف ثالثة محددة بشكل منفصل وفقًا للولاية القضائية المعمول بها، والقدرة على التأهل، والرسوم، والحدود، والشروط. العقود CLC App و Protocol v1.1.0 لا تعمل نفسها على السكة الحديدية.

يظل التنفيذ المتعدد المكالمات، والتأمين المشترك، وعلامة إدارة CLC المقترحة ومجموعة شبكة CLC المقترحة، وشبكة البطاقات، ومجموعات الدفع الصغيرة الشخصية، وخدمات الدفع المنظمة مقترحة أو تعتمد على النشر حتى يتم تحديد تنفيذ وشروطها الحاكمة بأنها نشطة.


\section*{16. نموذج التقييم المقترح}
\addcontentsline{toc}{section}{16. نموذج التقييم المقترح}

يمكن لبرنامج السجل أو المجموعة أو برنامج السيولة المستقبلي تكييف هذا الشكل. إنها ليست معايير أو ضمانات عالمية حالية للتسجيل.

\begin{enumerate}[leftmargin=*]
\item \textbf{الحقائق الملاحظة:} هوية الجهة المصدرة وتاريخها، شروط البطاقات، العروض، دليل القدرة، تشكيل المجموعة، المراقبين، المخزون، الحدود، الاحتياطيات، الضامن، الحوادث، والالتزامات غير الحل.
\item \textbf{الغرض والأطراف المتضررة:} الفائدة المقصودة والمخاطر والصراعات والآثار البيئية والاجتماعية، والذي يحمل الخسارة أو المسؤولية.
\item \textbf{التقييم:} سعر العرض، القيمة المعلنة على البطاقة، طريقة سعر الصرف في المجموعة، إشارة القرآن، الوحدات، العلامات الزمنية، وأسباب أي خلاف.
\item \textbf{الحدود:} التأهل، الاستخدامات المحظورة، التركيز، أقصى مستويات الميزان، التوقفات، القيود الجغرافية أو القانونية، وتحفيزات المراجعة.
\item \textbf{الحماية:} الحالة المحددة للجهة الملتزمة، الحدث المشمول، التمويل، الحد الأقصى، الإقصاءات، المدة، الأدلة، عملية المطالبة، والإجراءات التعريفة.
\item \textbf{القرار:} الموافقة أو مراجعة أو رفض أو توقف أو تصاعد للمراجعة البشرية. تحديد صانع القرار والأسباب والشروط وتاريخ المراجعة وعملية الاستئناف أو التعديل.
\end{enumerate}
لا ينبغي أن تصف الناتج الاحتياطي أو الحدود أو الاحتياطيات أو التحقق كإقرار أو قدرة الجهة المصدرة أو التأمين أو الوفاء المضمون ما لم يثبت الطرف المسؤول صراحة هذه الحماية.


\section*{17. شهادة القانون والامتثال}
\addcontentsline{toc}{section}{17. شهادة القانون والامتثال}

يمكن لـ CPP تنسيق الرموز، والبطاقات، والمساحات، والتبادلات التي يعتمد التعامل القانوني عليها على تصميمها والتسويق واستخدامها، والأطراف المسؤولة، والولاية القضائية. لا شيء في هذا الورقة البيضاء هو نصيحة قانونية أو ضريبية أو استثمارية أو ائتمانية أو مالية.

استخدام CLC App يتحكم في \href{https://docs.cosmolocal.credit/ar/governance/terms}{شروط الخدمة}. . تشغل GEF التطبيق والبنية التحتية الداعمة. ما لم تتحمل GEF صراحة دورًا آخر، فإنها ليست مصدرًا، مشرف الصندوق، أمينًا، مقرضًا، مقترضًا، وسيطًا، ضامنًا، فديًا، مستشارًا، مؤسسًا تأمينًا، أو طرفًا في التزامات بين المشاركين.

يجب على خدمات الدفع التي تعتمد على التنفيذ تحديد مقدمها المسؤول، والولايات القضائية، والقدرة على المشاركة، والرسوم، والحدود، ونموذج الاحتفاظ بها، والشروط. وجود اتصال لا يعني أن GEF أو مجموعة تعمل على الخدمة المنظمة الأساسية.

\subsection*{17.1 الإفصاحات عن البطاقات والعروض}
\addcontentsline{toc}{subsection}{17.1 الإفصاحات عن البطاقات والعروض}

يجب على الجهة المصدرة أن يعلن ويحافظ على مواصلة:

\begin{itemize}[leftmargin=*]
\item المعلومات المسؤولة عن الهوية والاتصال
\item العرض المرتبط، ووحدة، والإمدادات، والقيمة المعلنة، والقدرة؛
\item مواقع وإجراءات العرض؛
\item توقيت الوفاء والدليل؛
\item قواعد انتهاء الصلاحية والرسوم والضرائب والقيود والتغيير المادي؛
\item الشكاوى أو الاستبدال أو التعويضات أو التعويضات الأخرى؛ و
\item طريقة الإبراء التي تمنع إعادة الاستخدام بعد الإنجاز.
\end{itemize}
العقد الرمزي لا يضع هذه الشروط أو يثبت الأداء.

\subsection*{17.2 الإفصاحات عن المجموعة}
\addcontentsline{toc}{subsection}{17.2 الإفصاحات عن المجموعة}

يجب أن تنشر مشرف الصندوق:

\begin{itemize}[leftmargin=*]
\item هدف المجموعة وصانعي القرار المسؤولين عنها.
\item مالك \texttt{SwapPool}، ومدير الوكالة، ومراقبي الاعتماد، ومستفيدون من الرسوم؛
\item الأصول المقبولة وقواعد تعليق أو إزالة؛
\item أساليب التقييم، ومصادر الاقتباسات، والرسوم، والحد الأقصى لميزانية الرمز، والمخزون، وسلطات سحب المالكين؛
\item حقوق المساهمة والخروج
\item كل احتياطي، وضمان، ضامن، وسياسة التأمين، وقاعدة تخصيص الخسائر؛ و
\item الإدارة، الترقية، الطوارئ، الشكاوى، الهجرة، وإنهاء العمليات.
\end{itemize}
الإدراج في القائمة لا يشكل تأكيدًا أو تقييمًا أو تأمينًا أو ضمانًا من GEF.

\subsection*{17.3 الإجراءات والالتزامات}
\addcontentsline{toc}{subsection}{17.3 الإجراءات والالتزامات}

رجل عادي \textbf{تبادل المجموعة} البورصات الأصول المدعومة بموجب الاقتراحات والحدود المعاملة المعروضة. إدارة الأصول لا تجعلها قرضًا أو سدادًا أو استردادًا للمصدر أو تحقيقًا في العالم الحقيقي.

\textbf{تقديم القسيمة للاسترداد} يعيد أو يعرض وحدات البطاقة للمصدر. \textbf{الوفاء} هو الأداء الموعود للمصدر. \textbf{الإبراء} سجلات وحدات مكتملة بحيث لا يمكن إعادة استخدامها. التطبيق \textbf{«استرداد»} العمل يعد حاليا تحويل إلى صاحب الرمز؛ هذا التحويل وحده لا يثبت الوفاء.

التطبيق \textbf{«سحب القسيمة من التداول»} الإجراء هو إزالة الكتالوج العكسي من القائمة إنها لا تحرق الأرصدة، أو تلغي المطالبات، أو تنفيذ الالتزامات.

إن القرض المستقبلي أو منتج ائتماني آخر يتطلب تقديم شروط إضافية وشروط المعاملة بشكل منفصل قبل الاستخدام. لا يوجد إرسال عادي أو مساهمة أو إيداع مجموعة أو تبادل مجموعة أو تقديم أو استيفاء أو إفراز قرض بسبب اتجاه أو نوع الأصول.

\subsection*{17.4 الحماية والأدلة والقانون المحلي}
\addcontentsline{toc}{subsection}{17.4 الحماية والأدلة والقانون المحلي}

الحد، الاحتياطي، إدخال السجل، إدراج التطبيقات، أو صفقة بلوكتشين ليست بطبيعة الحال ضمان أو سياسة التأمين. يجب على أي حماية تحديد الجهة المسؤولة عنها، والحدث المشمول، والتمويل، والحد الأقصى، والإقصاءات، والمدة، والأدلة، وعملية المطالبة، والشروط المعمول بها.

الأدلة العامة السلسلة يمكن أن تظهر العناوين، والأصول، والكميات، والخطوط الزمنية، وأحداث العقد. فإنه لا يثبت نفسه هويته أو قدرة الجهة المصدرة أو الوفاء أو الرضا أو المناقشة الحاكمة أو الإبراء القانوني أو التأثير الاجتماعي.

قد تكون الميزات مقيدة أو غير متوفرة بسبب القانون أو العقوبات أو الأهلية أو حالة العقد أو المخزون أو الحدود أو الأمن أو الاختصاص القضائي أو خدمات الطرف الثالث. يظل الجهات المصدرة و مشرفو الصناديق والمقدمون والمشاركون مسؤولين عن تحديد وامتثال القانون المعمول به.


\section*{18. الخلاصة}
\addcontentsline{toc}{section}{18. الخلاصة}

يربط Cosmo-Local Credit التطبيق الحالي ومؤسسة Protocol v1.1.0 مع تصميم مقترح أوسع لحكم صناديق الالتزامات بشكل مستقل. يمكن أن تدعم المبادلات المباشرة الحالية للمجموعة، ومكونات الاقتباسات والحدود، وسجلات المعاملات العامة التبادل المسؤول، لكنها لا تثبت الوفاء بالجهة المصدرة، أو تخلق حقوق المساهمين، أو تضمن القيمة، والسيولة، والإفراج، والتأمين، والحالة القانونية، أو الحماية من الخسائر.

وتتطلب التوجيهات التنفيذية المقترحة، والتسويق، وأصول الحوكمة، وتصفية الشبكة، وبرامج السيولة، والحماية المشتركة الموصوفة في هذه الورقة تنفيذ منفصل، والتمويل، والحوكمة، والشروط المنشورة. يهدف التصميم إلى توسيع القدرة على التشغيل مع الحفاظ على أداء الجهة المصدرة، وإدارة المجموعة، والمساءلة في العالم الحقيقي مع الأطراف المحددة.


\section*{(أ) النموذج المقترح للقياس والرسوم}
\addcontentsline{toc}{section}{(أ) النموذج المقترح للقياس والرسوم}

يحدد هذا الملحق إطار قياس مقترح. لا تسجل Protocol v1.1.0 الوفاء بالجهة المصدرة، أو الإبراء في العالم الحقيقي، أو كل حقل بيانات مطلوبة أدناه.

\subsection*{تعريف الأحداث والأسهم}
\addcontentsline{toc}{subsection}{تعريف الأحداث والأسهم}

بالنسبة لفئة البطاقات \emph{j}، الجماعة أو الفترة \emph{t}، وطريقة التقييم المعلنة \emph{m}:

\begin{itemize}[leftmargin=*]
\item \texttt{O\_\{j,t,m\}}: قيمة الالتزامات المؤهلة المعلقة عند حد قياسها.
\item \texttt{X\_\{j,t,m\}}: قيمة عمليات تبادل المجموعة المكتملة خلال الفترة.
\item \texttt{P\_\{j,t\}}: الوحدات التي تقدم بصورة صحيحة للمصدر للاسترداد.
\item \texttt{F\_\{j,t\}}: الوحدات المقدمة مع إثبات الوفاء المنصدر بشكل منفصل.
\item \texttt{G\_\{j,t\}}: وحدات مكتملة ذات سجل الإفراز الذي يمنع إعادة الاستخدام.
\end{itemize}
لا يمكن استنتاج \texttt{O} من إمدادات الرموز وحدها. يجب أن تحدد سياسة قياس الجهة المصدرة المسؤول وتستبعد، حسب الاقتضاء، المخزون الذي يحتفظ به الجهة المصدرة، والحوافز المحروقة، والحوافز المنتهية من الصلاحية، والحوافز المفروضة، والحوافز الاختبارية، والحوافز غير الوصولية، والرموز التي لا تخلق شروطها التزامًا مستمرًا للطرف الثالث.

يجب أن تنشر كل قياس قيمة الوحدة ومصدرها وأسلوب تقييمها وطابع زمني ومعاملة أسعار الصرف المختلفة. يمكن لنقل على السلسلة دعم \texttt{X} أو تقديم الأدلة. فإنه لا يؤكد بذاته \texttt{F} أو \texttt{G}.

\subsection*{تدابير الوفاء القائمة على الجماعات}
\addcontentsline{toc}{subsection}{تدابير الوفاء القائمة على الجماعات}

بالنسبة لمجموعة من تقديمات الإفراج الصالحة:

\texttt{FulfillmentRate = FulfilledPresentments / ValidPresentments}

\texttt{DischargeCompleteness = DischargedFulfillments / FulfilledPresentments}

استخدم نفس المجموعة المغلقة أو الناضجة في كل عدد ومسمى. تقرير رفض، سحب، انتهت صلاحيته، تنازع، تم تنفيذها جزئيًا، تم تصحيحها، ولا تزال مفتوحة.

\texttt{FulfillmentLatency = median(t\_fulfilled - t\_presented)}

\texttt{HoldingDuration = median(t\_presented - t\_acquired)}

قياسات تأخر الوفاء خدمة الجهة المصدرة بعد العرض. مدة الاحتفاظ هي تدبير منفصل ولا يجوز وضع علامة على تأخر الاسترداد.

\subsection*{قياسات السرعة المختلفة}
\addcontentsline{toc}{subsection}{قياسات السرعة المختلفة}

لا يمكن حساب سرعة الإبراء المقترحة من الالتزام إلا عندما تستخدم \texttt{O} والقيمة المكتملة نفس طريقة التقييم:

\[
V_{\text{\latinfont commit},t} = \frac{\text{\latinfont FulfilledValue}_t}{\text{\latinfont AverageOutstandingEligibleCommitmentValue}_t}
\]

تدابير النشاط المتبادل في المجموعة منفصلة:

\texttt{V\_swap,t = PoolSwapValue\_t / AverageMeasuredPoolInventoryValue\_t}

ولا يثبت أي من هذه القيم تأثير اجتماعي أو قدرة الجهة المصدرة أو ربحية أو تحويل النقدية.

\subsection*{إيرادات الشبكة المقترحة}
\addcontentsline{toc}{subsection}{إيرادات الشبكة المقترحة}

أترك:

\begin{itemize}[leftmargin=*]
\item \texttt{PF\_t} هي رسوم المجموعة الإجمالية التي تم إنشاؤها خلال الفترة.
\item "\texttt{NR\_t}" هي المجموعة المقترحة للشبكة التي تم استلامها فعلياً كمشاركة من هذه الرسوم المكتشفة؛
\item \texttt{RF\_t} تكون الرسوم المقدمة للطرق أو الخدمات التي تم الحصول عليها بشكل منفصل؛ و
\item \texttt{$\chi$\_t} هو النسبة المحددة من الإيرادات المستلمة التي تكون مؤهلة وتتحول للاستخدام المعلن بالعملات النقدية بعد التكاليف والقيود السياسية.
\end{itemize}
من منظور ميزانية الشبكة المقترحة:

\texttt{ProposedNetworkRevenue\_t = NR\_t + RF\_t}

\texttt{CashUsableNetworkRevenue\_t = $\chi$\_t $\times$ (NR\_t + RF\_t)}

لا تضيف \texttt{PF\_t} إلى \texttt{NR\_t}: النقود هي تحويل من رسوم المجموعة الإجمالية وإلا سيتم احتسابها مرتين. بدلاً من ذلك ، تدعم Protocol v1.1.0 الحالية رسوم بروتوكول إضافية ؛ يجب الإبلاغ عن إيراداتها بشكل منفصل عن هذا النموذج المقترح.


\section*{ب. تساقط ضريبة مستقبلية مثالية}
\addcontentsline{toc}{section}{ب. تساقط ضريبة مستقبلية مثالية}

هذا هو التصميم المقترح لنشر مستقبلي لا تنطبق إلا إذا تم تنفيذها وتمويلها واتخاذها من خلال الحوكمة المنشورة وشروط المشاركين. فإنه لا يصف رسوم Protocol v1.1.0 الحالية، أو حق المساهمين، أو احتياطي، أو ضمان، أو ترتيب التأمين.

دعونا \texttt{F\_in} تكون الإيرادات التي تلقتها بالفعل ميزانية الشبكة المقترحة خلال فترة: الركوب المقترح للشبكة، رسوم التوجيه أو خدمات منفصلة، وأي إيرادات أخرى محددة صراحة. الرسوم الإجمالية للمجمعات التي تبقى مع المجمعات غير مدرجة.

يجب تصنيف أصول الإيرادات على أنها:

\begin{itemize}[leftmargin=*]
\item \texttt{E\_cash}: الأصول المؤهلة للاستخدام المعلن بالعملات النقدية بموجب السياسة المعتمدة أو
\item \texttt{E\_kind}: قسائم طبيعية أو أصول أخرى لا تعتبر قابلة للتحويل النقدي.
\end{itemize}
أي سياسة تحويل تحتاج إلى مسؤولي الأصول والموقع، والسلطات المسؤولة، ومصادر الأسعار، وحدود الانزلاق، والتقارير، والضوابط القانونية المطبقة.

قد يطبق الشلال المقترح الإيرادات المستلمة في هذا الترتيب:

\begin{enumerate}[leftmargin=*]
\item \textbf{هدف الاحتياطي المغطى:} تمويل الاحتياطي المعتمد أو هدف التأمين القائم فقط على التعرض المغطى المحدد والأصول المؤهلة والإقصاءات وقواعد المطالبات.
\item \textbf{العمليات الأساسية:} تمويل ميزانية تشغيل محددة ومعلنة.
\item \textbf{ولايات السيولة:} تمت الموافقة على البرامج المجمعية أو التوجيهة التي يتم تحكمها بشكل منفصل.
\item \textbf{الميزانية المتبقية:} تخصيص أي مبلغ متبقي بين العمليات وبرامج السيولة ومدفوعات إضافية تحت القيود المنشورة.
\end{enumerate}
هدف الاحتياطي ليس بحد ذاته ضماناً. أي تأمين أو ضمان يجب أن يحدد الطرف الملتزم، والحدث المغطى، والتمويل، والحد الأقصى، والإقصاءات، والمدة، والدليل، وعملية المطالبة، وتخصيص الخسائر.

\textbf{الحراسة:} تُخصص تخصيصات الشلالات لتغطية العمليات والخدمات السيولة المعتمدة. لا يجوز إطارها أو تنفيذها كعمليات دعم الأسعار.

بموجب النموذج المقترح، سيتم استقالة أي رموز إدارة CLC المقترحة التي تم الحصول عليها من خلال برنامج السيولة الخارجية المفعول بشكل منفصل أو وضعها في مخزن غير مصادر للتصويت. هذه الآلية لا تنشط


\section*{(ج) التعريفات المقترحة KPI}
\addcontentsline{toc}{section}{(ج) التعريفات المقترحة KPI}

هذه المؤشرات الرئيسية تشكل مواصفات قياس مقترحة، وليس بيانًا بأن التطبيق الحالي أو البروتوكول يسجل كل حدث مطلوب.

يجب أن تتضمن كل KPI المنشورة:

\begin{itemize}[leftmargin=*]
\item الناشر المسؤول ومصدر البيانات
\item تعريف الحدث أو الحالة
\item فترة المجموعة أو فترة القياس؛
\item الوحدة و طريقة التقييم
\item طابع زمني للتقييم
\item قواعد الإدماج والاستبعاد
\item معالجة السجلات الجزئية والمتنازع عليها والمتأخرة والمنع من الوصول إليها والمصالحة؛
\item دليل أو شهادة مطلوبة خارج السلسلة
\item تاريخ المراجعة وتقييدات جودة البيانات و
\item ما إذا كانت النتيجة موجودة على السلسلة أو تم الإبلاغ عنها أو إثباتها أو التحقق منها بشكل مستقل أو التقدير.
\end{itemize}
\begin{longtable}{P{0.31\linewidth} P{0.31\linewidth} P{0.31\linewidth}}
\toprule
KPI & التعريف المقترح & الأدلة المطلوبة والاستبعاد \\
\midrule
\endhead
\textbf{تقديمات صالحة} & الحساب أو الوحدات المقبولة في عملية استرداد الجهة المصدرة خلال فترة & تحديد العرض، والجهة المصدرة، وترخيص الاحتفاظ بها، والكمية، والوقت، والحالة. استبعاد النسخ المزدوجة والطلبات غير الصالحة \\
\textbf{معدل الوفاء} & التقديرات الصالحة المكتملة مقسمة على التقديرات الصالحة لنفس الفئة الناضجة & أدلة منفصلة على أداء الجهة المصدرة الإبلاغ عن القضايا المفتوحة والرفضية والمتنازع عليها والجزئية والصحيحة \\
\textbf{اكتمال الإبراء} & العروض المكتملة ذات سجل الإبراء مقسمة على العروض المكتملة & حرق، إلغاء، سجل إعاقة، أو دليل آخر على عدم إعادة الاستخدام مرتبط بالوفاء \\
\textbf{تأخير الإنجاز} & المتوسط و90 في المئة من الوقت من العرض إلى الوفاء & لا تستبدل مدة الحفاظ على الإصدار إلى التقديم أو الاستحواذ إلى التقديم \\
\textbf{مدة الاحتفاظ} & المتوسط وتوزيع الوقت من الاستحواذ إلى التقديم & تحديد حدث الاستحواذ واستبعاد أوقات الاستحواذ غير المعروفة \\
\textbf{الالتزامات المؤهلة المعلقة} & التزامات الطرف الثالث المؤهلة التي تبقى بعد إقصاءات محددة & هوية الجهة المصدرة وشروطه استبعاد المخزون المعمول به من الجهة المصدرة، وإنتهاء الصلاحية، الحرق، الإبراء، الاختبارات، ورمز عدم الالتزام. \\
\textbf{حجم تبادل المجموعة} & قيمة التبادلات المباشرة المكتملة في المجموعة بموجب طريقة تقييم واحدة تكشف عنها & أحداث التسوية على السلسلة، ووحدات الرمز، ومصدر السعر، والوقت. لا تصنفها كإنجاز للمصدر \\
\textbf{مخزون السباحة} & الأصول المدعومة القياسية التي تمتلكها مجموعة في وقت محدد & رصيد العقود، احتياطيات الرسوم، الأصول غير المتاحة، طريقة التقييم، وسلطات سحب المالكين \\
\textbf{كفاية الاحتياطي} & الأصول الاحتياطية المتاحة المؤهلة مقسمة على التعريضات المغطاة صراحة & سياسة التغطية، مؤهلية الأصول، الاحتفاظ/السيطرة، الالتزامات، الاستبعاد، وتقييم؛ ليس إجمالي إمدادات الرموز الافتراضية \\
\textbf{استخدام الحد} & ميزان رموز المجموعة المقاسة مقسومًا على السقف الذي تم تشكيله حاليًا & العنوان Limiter، الرمز، المجموعة، العلامة الزمنية، التغييرات، والفترات دون مقيد \\
\textbf{معدل عبور الاقتباس} & ردود الاقتباس الناجحة مقسمة على محاولات الاقتباس الصالحة & النتيجة التي تقتصر على الاقتباسات لا تنفيذ الطريق \\
\textbf{معدل تنفيذ الطريق} & التنفيذات المتعددة الهوب المكتملة مقسمة على محاولات التنفيذ الصحيحة & لا يمكن تطبيقها إلا على نظام تنفيذ تم تنفيذه. تقرير عن قواعد الحد الأدنى و القواعد الذرية \\
\textbf{استرداد الضامن} & تعويضات مؤهلة للحصول عليها مقسمة على المطالبات المغطاة المدفوعة & الضامن المحدد، سياسة المطالبات، التوقيت، التكاليف، النزاعات، وإزالة التكاليف \\
\textbf{إيرادات الشبكة المقترحة} & المقترحات المتعلقة بالشبكة التي تم استلامها بالإضافة إلى رسوم التوجيه / الخدمة المنفصلة المستلمة & استبعاد رسوم البحيرة الإجمالية التي تم احتجازها من قبل البحيرات وتجنب احتساب الركبة مرتين \\
\textbf{توقيت الإدارة} & وقت الكشف عن الحادث، القرار، التوقف، الإصلاح، وإغلاقه & ساعات محددة، الهيئات المسؤولة، سلطات الطوارئ، الاستئناف، والأحداث المفقودة \\
\bottomrule
\end{longtable}

مطالبات التأثير الاجتماعي تتطلب منهجية منفصلة. أنشطة بلوكتشين وحدها لا تثبت الهوية أو أداء الجهة المصدرة أو الرضا أو صحة المجتمع أو السببية أو التأثير.


\section*{(د) المعايير المقترحة للإطلاق}
\addcontentsline{toc}{section}{(د) المعايير المقترحة للإطلاق}

كل قيمة أدناه تشير إلى ذلك إنها ليست تطبيقًا افتراضيًا حاليًا ، إعداد Protocol v1.1.0 ، وضمان ، عرض ، أو التزام التسليم. سيتعين على النشر في المستقبل تبني وتنفيذ ومراقبة ونشر معاييرها الفعلية.

\begin{itemize}[leftmargin=*]
\item \textbf{مستويات الكواروم لعلامة الحكم المقترحة CLC:}
\begin{itemize}[leftmargin=*]
\item روتين Q1: على الأقل 4 في المائة من القرارات وأكثر من 50 في المائة من الموافقة.
\item حساسة Q2: حكم أقصى بنسبة 10 ٪ وموافقة أقصى بنسبة 60 ٪.
\item حرجة Q3: موافقة بنسبة 20٪ على الأقل والقرار بنسبة 66.7٪ على الأقل.
\end{itemize}
\item \textbf{الأوقات المقترحة:}
\begin{itemize}[leftmargin=*]
\item T1: 48 ساعة لأعمال Q1.
\item T2: 7 أيام بالنسبة لأعمال Q2.
\item T3: 30 يومًا لفعاليات Q3
\end{itemize}
\item \textbf{تقدم العصر:} 7 الأيام، بما في ذلك التصويت المعتمد، ونشر الميزانية، والنوافذ المقترحة sCLC.
\item \textbf{توقف طارئ:} على الفور من خلال سلطة طوارئ محددة تنتهي صلاحيتها بعد 72 ساعة ما لم يتم التصديق عليها بموجب العملية المعتمدة.
\item \textbf{المقترح لتحقيق الشبكة:} 20 ٪ من رسوم المجموعة المشاركة بطبيعة الحال، محدودة بالسياسة.
\item \textbf{نطاق رسوم المجموعة التوضيحية:} 0\%-20\%، التي نشرتها كل مجموعة مشاركة.
\item \textbf{الحد الأقصى المقترح من رسوم التوجيه أو الخدمة:} 20 ب.ب.ب.س في كل طريق
\item \textbf{المعدل المقترح لسرعة الشبكة:} \texttt{rake\_rate\_p = pool\_fee\_p $\times$ rake\_share\_p}.
\item \textbf{مؤهلات الإيرادات:} تصنيف الأصول المستلمة كـ \texttt{E\_cash} المؤهلة للحصول على النقود أو \texttt{E\_kind} في الطبيعة حسب الطريقة المنشورة.
\item \textbf{التحكم في التحويل:} مسؤولي الأصول والموقع، والسلطات المسؤولة، ومصادر الأسعار، ونوافذ الزمن، وحدود الانزلاق، القيود، والإبلاغ.
\item \textbf{تمكين الميزانية:} تنشر الميزانية المقترحة للحصول على الرسوم فقط بعد استيفاء الاحتياطي المغطى والمستهداف التشغيلية المعتمدة. الميزانية قد تكون صفر
\item \textbf{مسارات المخاطر:} لا تعرض فئة أصول واحدة لخطر فئة أخرى دون اختيار صريح ومعرفة بموجب القانون المعمول به.
\item \textbf{القيود المقترحة للتداول والحساب:} يرفضون أنشطة متعددة الفئات غير المعتمدة ويطبقون القيود المعتمدة.
\item \textbf{خفض التغطية الاختياري:} فقط عندما تسمح بها شروط التغطية القائمة مسبقاً والاتفاق المطلوب والقانون المعمول به؛ لا يقلل بمفرده من الالتزام بالبطاقة أو الرصيد على السلسلة.
\item \textbf{مراقبة السيولة الخارجية، إذا تم تمكينها بشكل منفصل:} نشر نطاق الموقع، طريقة القياس، المحفزات، القيود المفروضة على الإنفاق والحجم، ومراقبة التنفيذ، والوقوفات الطارئة، والإيصالات. لا تعرض البرنامج كدعم رمزي للثمن.
\end{itemize}
تبقى رسوم Protocol v1.1.0 الجمعية الحالية، رسوم البروتوكول الإضافية، والحد الأقصى لموازنة رموز الجمعية تحكمها العقود والتنظيمات المنشأة، وليس هذه القيم المقترحة.


\section*{(إي) نموذج عمل --- مقترح ريك الشبكة}
\addcontentsline{toc}{section}{(إي) نموذج عمل --- مقترح ريك الشبكة}

هذا المثال يوضح نموذج الرف المقترح. هذه ليست الحسابات الإضافية للبروتوكول المخصصة التي تستخدمها Protocol v1.1.0.

افترض:

\begin{itemize}[leftmargin=*]
\item المجموعة المشاركة تتقاضى رسوم 2.00\% المجموعة؛
\item يتلقى الرقيب المقترح للشبكة 20٪ من رسوم المجموعة؛ و
\item 25 في المائة من النقود المستلمة هي مؤهلة و قابلة للتحويل للاستخدام المعلن بالنقود بعد التكاليف.
\end{itemize}
المعدل الفعلي المقترح لسرعة الشبكة على القيمة المتوجهة هو:

\texttt{rake\_rate = 2.00\% $\times$ 20\% = 0.40\% = 40 bps}

الجزء القابل للاستخدام النقدي تحت عامل الطلب المفترض بنسبة 25٪ هو:

\texttt{cash\_usable\_rate = 40 bps $\times$ 25\% = 10 bps}

إذا كان لدى ميزانية الشبكة المقترحة متطلبات نقدية شهرية تبلغ 150،000 دولار ولا توجد إيرادات أخرى، فإن القيمة التوجيهية التوضيحية المطلوبة بمعدل قابل للاستخدام النقدي بنحو 10 بيبس هي:

\texttt{required\_routed\_value = \$150,000 / 0.001 = \$150,000,000 per month}

هذا الحساب لا يشمل رسوم المجموعة الإجمالية التي تحتفظ بها المجموعات. إضافة هذه الرسوم إلى الشبكة من شأنها أن تضاعف حساب المصدر.

من الضروري أيضاً نشر تحليل ميزاني حقيقي:

\begin{itemize}[leftmargin=*]
\item الإيصالات الفعلية عن الركوب والرسوم الخدمية؛
\item مؤهلية الأصول وتكاليف التحويل
\item مشاركة المجموعة ونطاق المسار
\item تم تبني أهداف احتياطي وتشغيل؛
\item الخسائر والنزاعات والتصحيحات والأصول غير المتاحة. و
\item سيناريوهات الحساسية بدلاً من النمو الموعود أو العائدات.
\end{itemize}
أي ميزانية مقترحة للحصول على الرسوم أو لبرنامج السيولة ستبقى أسفل من أولوياتها المعتمدة ويمكن أن تكون صفر.


\section*{(ف) قائمة التحقق من غرفة البيانات}
\addcontentsline{toc}{section}{(ف) قائمة التحقق من غرفة البيانات}

\begin{enumerate}[leftmargin=*]
\item \textbf{دليل تاريخي Sarafu Network}
\begin{itemize}[leftmargin=*]
\item مبلغ المقايضات، وصف المستخدم النشط والأصول، الحوادث، تقديمات الاسترداد، قسائم الشيخوخة، الإنجاز المؤكد، الإبراء، تدابير مدة الاحتفاظ، الفترات، الاستبعاد، والمنهج بشكل منفصل عن المقاييس Cosmo-Local Credit الحالية.
\item الحفاظ على \href{https://dune.com/grassrootseconomics/sarafu-network}{لوحة التحكم التاريخية Dune Analytics}.
\end{itemize}
\item \textbf{الأدلة المقترحة لتنفيذ الرسوم، إذا تم تنفيذها}
\begin{itemize}[leftmargin=*]
\item إظهار أن أي خط الرسوم والبوابات السجلية متكاملة في مسار المعاملات التي يتم تنفيذها بدلاً من فرضها فقط بواسطة واجهة.
\item تقديم اختبارات تظهر أن الخدمة التنفيذية المعمول بها ترفض الطرد الذي لا يثبت استلامه جمع الرسوم بموجب السياسة المعتمدة.
\item نشر متجهات الاختبار، حالات الفشل، الأصول المؤهلة، ومجمعات متوافقة مرتبطة بنسخة السجل المحددة.
\item ربط التطبيقات \href{https://software.grassecon.org}{وصف البرمجيات}.
\end{itemize}
\item \textbf{الضمان والحزمة التغطية}
\begin{itemize}[leftmargin=*]
\item نشر المؤهلات، والأطراف المسؤولة، والتمويل، والحد الأقصى، والإجازات عند الاقتضاء، والتحفيزات، والأدلة، والاستبعادات، سلطة القرار، ومثائل النزاعات والإستئناف.
\item تشمل قسيمة العينة والكشف عن المجموعة التي تميز مسؤولية الجهة المصدرة عن أي حماية المجموعة أو ضمان من طرف ثالث يتم تقديمها صراحة.
\item الحفاظ على \href{https://sarafu.network/reports}{التاريخ Sarafu Network تقارير أرشيف} و \href{https://youtu.be/5yGaP31bvs0?si=fZ-AMTEXsmVR8Ulv}{تقديم السنة التاريخية في المراجعة} كدليل على السلالة، وليس دليل على تغطية CLC الحالية أو التبني.
\end{itemize}
\item \textbf{حزمة القانون والامتثال}
\begin{itemize}[leftmargin=*]
\item استراتيجية الولاية القضائية، فئات البطاقات، سياسات المكافئة النقدية، KYC أو خيارات التصديق، الجغرافية، الإفصاح عن المستهلكين، والسيطرة على البيع الخاطئ، والتمييز بين تبادل المجموعة العادية ومنتج الائتمان الصريح.
\item ربط الفعالية \href{https://docs.cosmolocal.credit/ar/governance/terms}{Cosmo-Local Credit شروط الخدمة} والحفاظ على الإصدارات المستبدلة من خلال السجلات المراقبة.
\end{itemize}
\item \textbf{تعزيز الحكم}
\begin{itemize}[leftmargin=*]
\item توفير إجراءات تضارب المصالح والرفض، سجلات القرارات والإستئنافات، والعقوبات، والإجراءات المناسبة لإزالة السجلات، ومحدود قوة الطوارئ، ومثال على تغييرات سعر الصرف والحد في الوقت المحدد.
\end{itemize}
\item \textbf{المصدر وحزمة التأمين}
\begin{itemize}[leftmargin=*]
\item نشر مخزون المصدر والترخيص والعناوين والإصدارات المتوفرة وتقارير المراجعة المتاحة والمتغيرات ومصدر البناء وسلطات الإدارة ومواصلات الحوادث وألواح المراقبة.
\item ربط التطبيقات \href{https://github.com/grassrootseconomics}{مستودعات الاقتصاد الجذري}.
\end{itemize}
\end{enumerate}

\end{document}
